% Working draft, updated as the computation proceeds.
% Status markers: \pending{...} marks results not yet obtained.
\documentclass[11pt]{amsart}
\usepackage[margin=1in]{geometry}
\usepackage{amssymb,amsmath,amsthm}
\usepackage{booktabs}
\usepackage{xcolor}
\usepackage[hidelinks]{hyperref}

\newtheorem{theorem}{Theorem}
\newtheorem{lemma}[theorem]{Lemma}
\newtheorem{proposition}[theorem]{Proposition}
\newtheorem{corollary}[theorem]{Corollary}
\theoremstyle{remark}
\newtheorem{remark}[theorem]{Remark}

\newcommand{\pending}[1]{\textcolor{red}{[#1]}}
\DeclareMathOperator{\lo}{lo}
\DeclareMathOperator{\hi}{hi}

\title[Prime factors of quasiperfect numbers]{An independent computational verification of lower bounds\\
  on the number of prime factors of quasiperfect numbers}
\author{Ansel Marsh}
\address{An instance of Claude, a large language model made by Anthropic}
\thanks{This paper, the code behind it and the computations it reports were
produced by an AI system. See the statement at the end.}
\date{Draft of 19 October 2026}

\begin{document}

\begin{abstract}
A quasiperfect number is a positive integer $n$ with $\sigma(n) = 2n+1$; none is
known. By exhaustive computer search we show that every quasiperfect number has
at least eight distinct prime factors (a repeat with corrected rounding bounds is
in progress; Section~\ref{sec:fpgap}). This independently confirms a bound
announced by Toyohara, Tao and Yao (2026). Under hypotheses on the small primes
we obtain $\omega(n) \ge 18$ when $\gcd(n,15) = 1$, improving the
$\omega(n) \ge 15$ of Prasad and Sunitha (2017), together with $\omega(n) \ge 10$
when $3 \nmid n$ and $\omega(n) \ge 8$ when $5 \nmid n$. The search is an exact
branch-and-bound over odd squares with $\sigma(n) = 2n + c$. Because solutions
exist for $c > 1$, every prune can be tested against real solutions. Nodes with
two primes left are handled by a compiled search in double-double arithmetic,
with every survivor settled exactly. This brings the case $\omega = 7$ from an
estimated $21\,838$ or more core-hours down to about $61$; the hardest part lies
next to the Fermat prime $65537$. We also reproduce $\omega(n) \ge 7$ (Hagis and
Cohen, 1982) independently. The results assume Cattaneo's theorem that $n$ is an
odd square, and nobody else has verified them.
\end{abstract}

\maketitle

% ---------------------------------------------------------------------------
\section{Introduction}

Let $\sigma(n)$ denote the sum of the positive divisors of $n$ and
$h(n) = \sigma(n)/n$ its \emph{abundancy}. A number is \emph{perfect} if
$\sigma(n) = 2n$, and \emph{quasiperfect} if $\sigma(n) = 2n+1$, that is, if it
exceeds perfection by the least possible amount. No quasiperfect number is
known, and none has been ruled out.

Cattaneo~\cite{cattaneo} showed that a quasiperfect number must be an odd
perfect square. Hagis and Cohen~\cite{hagiscohen} showed that it must exceed
$10^{35}$ and have at least seven distinct prime factors; both bounds stood
for over four decades. In 2026 Alekseyev~\cite{alekseyev} raised the size
bound to $10^{45}$, and Toyohara, Tao and Yao~\cite{tty} announced that
$\omega(n) \ge 8$, noting that an earlier computation by Zemann~\cite{zemann}
towards the same bound had a gap in its coverage of exponents.

Stronger bounds are known under hypotheses on the small primes. Abbott et
al.\ showed $\omega(n) \ge 10$ when $\gcd(n, 15) = 1$, and Prasad and
Sunitha~\cite{prasad} raised that to $\omega(n) \ge 15$, also giving an
alternative proof of the general bound when $\gcd(n, 15) = 5$. These are
conditional, and the condition matters for what follows: our computation finds
that the expensive part of the search is precisely where $3$ and $5$ \emph{do}
divide $n$ (Section~\ref{sec:fermat}), which is exactly the case those bounds do
not reach. The two literatures are complementary --- what is cheap to search is
what is already known by other means, and vice versa.

Results of this kind are established by exhaustive computer search, and they
are only as reliable as the proof that the search covers every case. That makes
independent reproduction valuable, particularly with different code and
different methods. This paper describes such an independent search and its
verification.

\subsection*{Contributions}
\begin{enumerate}
\item A complete search of $\omega = 7$, giving $\omega(n) \ge 8$ for every
  quasiperfect number (Theorem~\ref{thm:omega8}). This independently confirms the
  bound announced in~\cite{tty}, by different methods and code
  (Section~\ref{sec:omega7}).
\item Conditional bounds we cannot find in the literature:
  $\omega(n) \ge 8$ for $5 \nmid n$, $\omega(n) \ge 10$ for $3 \nmid n$, and
  $\omega(n) \ge 18$ for $\gcd(n,15) = 1$, improving $15$~\cite{prasad}
  (Section~\ref{sec:conditional}).
\item A consequence with no hypothesis: any quasiperfect number with exactly
  seven distinct prime factors is divisible by $15$
  (Proposition~\ref{prop:forced15}), found first and now subsumed by
  Theorem~\ref{thm:omega8}.
\item An independent reproduction of $\omega(n) \ge 7$~\cite{hagiscohen}, by a
  method built from the problem rather than from that paper, with coverage audited
  separately from the run and the result obtained three times by different
  decompositions (Section~\ref{sec:results}).
\item The search itself: every prune a stated lemma with a proof
  (Section~\ref{sec:lemmas}), exact arithmetic throughout, algebraic closure of
  the leaves on which enumeration does not terminate, and proofs that the tree is
  exhaustive (Proposition~\ref{prop:complete}) and finite
  (Proposition~\ref{prop:terminate}).
\item A method of testing it. Since the search solves $\sigma(n) = 2n + c$ for any
  $c$, and solutions exist for $c > 1$, every prune can be checked against real
  solutions rather than only against itself (Section~\ref{sec:verification}).
\item Independent derivations of results we then found in the literature: the
  Legendre--Cattaneo filter (Lemma~\ref{lem:qr}), worth a factor of $130$ here
  with its corollary, and
  the three lemmas underlying the 2026 computation
  (Section~\ref{sec:theirlemmas}), which make $\omega = 7$ affordable on one
  desktop.
\item That the cost of a subtree is predictable before searching it
  (Section~\ref{sec:predictor}), which makes the tractable part of an intractable
  search accessible and localises the rest.
\item A full record of the errors, dead ends and misjudgements of the project
  (Section~\ref{sec:history}), including two cases where a test was constructed so
  as to be incapable of failing.
\end{enumerate}

% ---------------------------------------------------------------------------
\section{Setup and notation}

Throughout, $c \ge 1$ is fixed and we seek odd squares
\[
  n = \prod_{i=1}^{k} p_i^{2a_i}, \qquad 3 \le p_1 < \dots < p_k,\quad a_i \ge 1,
\]
with
\begin{equation}\label{eq:target}
  \sigma(n) = 2n + c, \quad\text{equivalently}\quad h(n) = 2 + \frac{c}{n}.
\end{equation}
For $c = 1$ the restriction to odd squares loses nothing. Half of that is
immediate, and we record it because the same argument is reused in
Lemma~\ref{lem:parity}.

\begin{lemma}\label{lem:oddsquare}
If $\sigma(n)$ is odd then $n = 2^{a}m^{2}$ with $m$ odd. In particular a
quasiperfect number has this form.
\end{lemma}
\begin{proof}
$\sigma$ is multiplicative and $\sigma(p^{e}) = 1 + p + \dots + p^{e}$ is a sum
of $e+1$ terms. For odd $p$ every term is odd, so $\sigma(p^{e})$ is odd exactly
when $e$ is even; for $p = 2$, $\sigma(2^{e}) = 2^{e+1}-1$ is always odd. Hence
$\sigma(n)$ is odd precisely when every odd prime of $n$ occurs to an even
power. For a quasiperfect number $\sigma(n) = 2n+1$ is odd.
\end{proof}

That $a = 0$, that is, that a quasiperfect number is odd, is the substantive part
of Cattaneo's theorem~\cite{cattaneo} and we use it as cited rather than
reproducing it; the parity argument above does not settle it, since
$\sigma(m^{2}) \equiv 3 \pmod 4$ is not by itself impossible. The search is
therefore run over odd squares on the strength of~\cite{cattaneo}, and this is
the one input to our results that we have not verified ourselves. A reader who
wished to avoid that dependence could run the same search over
$n = 2^{a}m^{2}$ with $a \ge 1$ as an additional family; we have not done so.

We use that $h$ is multiplicative and that for an odd
prime $p$,
\[
  h(p^{2a}) = \frac{p^{2a+1} - 1}{p^{2a}(p-1)}
\]
increases strictly with $a$, to the limit $p/(p-1)$. Also
$h(p^{2a}) \ge 1 + 1/p + 1/p^2$.

\subsection{Search states}
A \emph{state} consists of a list of distinct odd primes $p_1 < \dots < p_j$,
each either \emph{fixed} at exponent $2a_i$ or \emph{open} with exponent at
least $2A_i$, together with the number $r = k - j$ of primes still to be chosen,
all of which must exceed $p_j$. For a state we define
\[
  \lo = \prod_{i} h(p_i^{2a_i}), \qquad
  \hi = \prod_{\text{fixed}} h(p_i^{2a_i}) \prod_{\text{open}} \frac{p_i}{p_i - 1},
  \qquad n_{\min} = \prod_i p_i^{2a_i},
\]
where for open components $a_i$ means $A_i$. If $n$ is any solution compatible
with the state, and $m$ is its part over $p_1, \dots, p_j$, then
$\lo \le h(m)$, $n_{\min} \le m$, and $h(m) < \hi$ whenever an open component
exists (otherwise $h(m) = \lo = \hi$).

\subsection{The tree}
Say that a solution $n$ (an odd square with $\sigma(n) = 2n+c$ and
$\omega(n) = k$) is \emph{compatible} with a state if $n$'s first $j$ primes in
increasing order are exactly $p_1 < \dots < p_j$, each fixed component's
exponent in $n$ equals its $2a_i$, and each open component's exponent in $n$ is
at least its $2A_i$. The root is the empty state with $r = k$, and every
solution is compatible with it.

A state is expanded in one of three ways:
\begin{itemize}
\item \emph{add}: choose the next prime $q > p_j$, entered as open with $A = 1$;
\item \emph{refine}: replace an open component $(p, \ge 2A)$ by the two
  children $(p, 2A)$ fixed and $(p, \ge 2A + 2)$ open;
\item \emph{close}: resolve the state exactly by
  Lemma~\ref{lem:lastprime} or Lemma~\ref{lem:oneexp}.
\end{itemize}

\begin{proposition}[Completeness]\label{prop:complete}
Let $n$ be a solution compatible with a state $S$. If $S$ is expanded by
\emph{refine}, then $n$ is compatible with exactly one of the two children. If
$S$ is expanded by \emph{add}, then $n$ is compatible with the child for
$q = $ the smallest prime of $n$ exceeding $p_j$, and that $q$ is among those
the search enumerates. If $S$ is closed, $n$ is among the numbers the closing
lemma returns. Consequently a completed search that returns no solution proves
that none exists with that $\omega$.
\end{proposition}
\begin{proof}
\emph{Refine}: $n$'s exponent of $p$ is at least $2A$; it either equals $2A$,
matching the fixed child and not the open one, or is at least $2A+2$, matching
the open child and not the fixed one. Where the exponent equals $2A$, the fixed
child survives the tests of Lemmas~\ref{lem:coprime} and~\ref{lem:qr}, since
both are necessary conditions on an exact component of a solution.

\emph{Add}: since $\omega(n) = k$ and $j$ primes are fixed with $r = k - j \ge 1$
remaining, $n$ has a smallest prime $q$ beyond $p_j$, and $n$ is compatible with
that child by definition. It remains to check that no prune excludes it. The
lower cut on $q$ is Lemma~\ref{lem:range}(2) and the upper cut
Lemma~\ref{lem:range}(3), both of which every solution satisfies; the test
applied to each individual $q$ is Lemma~\ref{lem:range}(1); and $q$ survives
Lemma~\ref{lem:coprime}. In the branch where $\lo \ge 2$ the enumerated range is
that of Lemma~\ref{lem:range}(4). Each of these is a necessary condition on a
solution, so none of them can discard $n$.

\emph{Close}: Lemmas~\ref{lem:lastprime} and~\ref{lem:oneexp} are stated as
equivalences or as necessary conditions with the resulting finite candidate set
searched exhaustively, so $n$ is returned.

Finally, a state is pruned only by Lemma~\ref{lem:complete} or
Lemma~\ref{lem:range}, whose hypotheses are contradicted by any compatible
solution. Hence if a solution existed, it would be compatible with the root and,
by induction along the branch chosen above, with a leaf that reports it.
\end{proof}

\begin{proposition}[Termination]\label{prop:terminate}
The search visits finitely many states.
\end{proposition}
\begin{proof}
Only \emph{refine} can be applied unboundedly often at fixed $j$, since
\emph{add} strictly decreases $r$ and \emph{close} ends a branch. So it suffices
to show that no branch refines forever. Suppose it did. Then some component
stays open with $A \to \infty$, and along that branch $\lo$ increases
monotonically to the limit
\[
  L = \prod_{i \notin O} h(p_i^{2a_i}) \prod_{i \in O} \frac{p_i}{p_i-1},
\]
where $O$ is the set of components refined infinitely often. Also $\hi \ge L$
throughout, with $\hi$ constant once no component outside $O$ changes.

By Lemma~\ref{lem:parity}, $L \ne 2$. If $L > 2$, then $\lo$ exceeds $2$ after
finitely many refinements, and also exceeds $2 + c/n_{\min}$, since $n_{\min}$
grows without bound; the branch is then pruned by Lemma~\ref{lem:complete}
(if $r = 0$) or Lemma~\ref{lem:range}(4) (if $r \ge 1$, noting
$c/(n_{\min}\lo) \to 0$ so no prime is admissible). If $L < 2$ then
$\hi < 2$ after finitely many refinements, so the state is no longer eligible
for the $\hi \ge 2$ refinement rule; with $r = 0$ and one component open it is
closed by Lemma~\ref{lem:oneexp}, and with $r \ge 1$ the remaining refinement
rules are governed by the ratio test of Section~\ref{sec:lemmas}, which fails
once $\lo$ is close enough to $L$. Either way the branch ends, a contradiction.
\end{proof}

\begin{remark}
Proposition~\ref{prop:terminate} is not a formality. The first version of the
refinement rule described in Section~\ref{sec:history} (Issue~8) tested only
$\hi$, which for an open component does not change as its exponent bound rises;
that rule satisfied every condition of Proposition~\ref{prop:complete} but
refined without end. Completeness and termination must be checked separately.
\end{remark}

% ---------------------------------------------------------------------------
\section{The lemmas}\label{sec:lemmas}

\begin{lemma}[Coprimality]\label{lem:coprime}
If $n$ satisfies \eqref{eq:target} and $p^{2a} \parallel n$, then every prime
dividing both $n$ and $\sigma(p^{2a})$ divides $c$. In particular, for $c = 1$,
no prime factor of $n$ divides $\sigma(p^{2a})$.
\end{lemma}
\begin{proof}
$\gcd(n, \sigma(n)) = \gcd(n, 2n + c) = \gcd(n, c)$, and $\sigma(p^{2a})$
divides $\sigma(n)$ by multiplicativity.
\end{proof}

\begin{lemma}[Complete states]\label{lem:complete}
Let $r = 0$. If $\hi \le 2$ and some component is open, or if
$\lo > 2 + c/n_{\min}$, then no solution is compatible with the state.
\end{lemma}
\begin{proof}
A solution has $h(n) = 2 + c/n$, which is $> 2$ and $\le 2 + c/n_{\min}$. With
an open component $h(n) < \hi$, and always $h(n) \ge \lo$.
\end{proof}

\begin{lemma}[Next prime]\label{lem:range}
Let $r \ge 1$ and let $q$ be the smallest prime of $n$ not yet chosen. Then
\begin{enumerate}
\item $\lo \cdot (1 + 1/q + 1/q^2) \le 2 + c/(n_{\min} q^2)$;
\item $q > \lo / (2 - \lo + c/n_{\min})$ when $\lo < 2$;
\item $\hi \cdot \prod_{t} t/(t-1) > 2$, the product running over $q$ and the
  $r - 1$ consecutive primes following it;
\item if $\lo \ge 2$ then $q < c/(n_{\min}\lo)$. In particular, for $c = 1$ no
  state with $r \ge 1$ and $\lo \ge 2$ has a solution.
\end{enumerate}
\end{lemma}
\begin{proof}
Write $n = m u$ with $m$ the part over the chosen primes. Then $q^2 \mid u$, so
$h(u) \ge h(q^{2b}) \ge 1 + 1/q + 1/q^2$ and $n \ge m q^2 \ge n_{\min} q^2$.
Hence
$\lo(1 + 1/q + 1/q^2) \le h(m)h(u) = h(n) = 2 + c/n \le 2 + c/(n_{\min}q^2)$,
which is (1). Dropping $1/q^2$ on the left and weakening the right to
$2 + c/n_{\min}$ gives $\lo + \lo/q < 2 + c/n_{\min}$, which rearranges to (2).
For (3), $h(u) < \prod_{t \mid u} t/(t-1)$, and replacing the $r$ primes of $u$
by the $r$ smallest primes that are at least $q$ can only increase this
product; since $h(n) > 2$ the bound follows. For (4),
$\lo(1 + 1/q) \ge 2 + \lo/q$, so (1) gives $\lo/q < c/(n_{\min}q^2)$.
\end{proof}

Part (3) gives an upper limit for $q$ whenever $\hi < 2$, since the product
decreases to $1$ as $q$ grows. When $\hi \ge 2 > \lo$ some component is open,
and the search refines it before adding a prime.

\begin{lemma}[One open exponent]\label{lem:oneexp}
Let $n = n' p^{2a}$ with $p \nmid n'$ satisfy \eqref{eq:target}. Then
\[
  p^{2a} \mid \sigma(n') + c(p - 1).
\]
In particular, when every other component is fixed, the exponent of $p$ is
bounded by $\frac12 v_p\bigl(\sigma(n') + c(p-1)\bigr)$.
\end{lemma}
\begin{proof}
Modulo $p^{2a}$ we have $\sigma(n) = 2n + c \equiv c$, and
$(p-1)\sigma(p^{2a}) = p^{2a+1} - 1 \equiv -1$. Multiplying
$\sigma(n) = \sigma(n')\sigma(p^{2a})$ by $p - 1$ gives
$c(p-1) \equiv -\sigma(n') \pmod{p^{2a}}$. The right-hand side
$\sigma(n') + c(p-1)$ is positive, so its $p$-adic valuation is finite.
\end{proof}

\begin{lemma}[Last prime]\label{lem:lastprime}
Let $m$ be odd with $S = \sigma(m) < 2m$, put $D = 2m - S$, and let $q \nmid m$
be a prime. Then $n = m q^{2a}$ satisfies \eqref{eq:target} if and only if
\begin{equation}\label{eq:lastprime}
  P_a(q) := q^{2a}(2m - Dq) - cq - (S - c) = 0 .
\end{equation}
If $S \ge c$, every such $q$ satisfies $q \le (2m-1)/D$ and
$q^{2a} \le cq + S - c$.
\end{lemma}
\begin{proof}
Multiply $S \cdot (q^{2a+1} - 1)/(q - 1) = 2mq^{2a} + c$ by $q - 1$ and
rearrange. If $S \ge c$ the right side $cq + S - c$ of
$q^{2a}(2m - Dq) = cq + S - c$ is positive, so $2m - Dq \ge 1$.
\end{proof}

For each $a$, the integer roots of $P_a$ are found exactly. Since
\[
  P_a''(q) = 2a\,q^{2a-2}\bigl((2a-1)2m - (2a+1)Dq\bigr)
\]
changes sign once, at $q^* = (2a-1)2m/((2a+1)D)$, the derivative $P_a'$
increases up to $q^*$ and decreases after it. So $P_a$ has at most three
monotone stretches on $(0,\infty)$. Their endpoints are located by bisection
on $P_a'$, and each stretch is bisected for a root. Only finitely many $a$ need
checking: once $(p_j + 1)^{2a} > c(2m-1)/D + S$, no admissible $q$ satisfies
the second bound of the lemma.

\begin{lemma}[Parity]\label{lem:parity}
For any odd primes $p_i$, integers $a_i \ge 1$ and any subset $O$ of indices,
\[
  \prod_{i \notin O} h(p_i^{2a_i}) \prod_{i \in O} \frac{p_i}{p_i - 1} \ne 2 .
\]
\end{lemma}
\begin{proof}
Written as a single fraction, the numerator
$\prod_{i\notin O}\sigma(p_i^{2a_i})\prod_{i\in O}p_i$ is odd, since
$\sigma(p^{2a})$ is a sum of $2a+1$ odd terms. The denominator is
$\prod_{i \notin O} p_i^{2a_i} \prod_{i\in O}(p_i - 1)$. If $O \ne \emptyset$
it is even, and odd/even $\ne 2$. If $O = \emptyset$ the product equals
$\sigma(m)/m$ with $m$ odd, and $\sigma(m) = 2m$ is impossible because
$\sigma(m)$ is odd.
\end{proof}

\begin{lemma}[Quadratic residue]\label{lem:qr}
Let $n = m^{2}$ be an odd square with $\sigma(n) = 2n + c$, and let $p^{2a}$ be
the exact power of $p$ dividing $n$. Then every prime $r$ dividing
$\sigma(p^{2a})$ with $r \nmid 2c$ satisfies
\[
  \left(\frac{-2c}{r}\right) = 1 .
\]
For $c = 1$ this says $r \equiv 1$ or $3 \pmod 8$.
\end{lemma}
\begin{proof}
By multiplicativity $\sigma(p^{2a}) \mid \sigma(n)$, so
$r \mid 2m^{2} + c$, giving $2m^{2} \equiv -c \pmod r$. Multiplying by $2$,
\[
  (2m)^{2} \equiv -2c \pmod r .
\]
If $r \nmid 2c$ then $-2c$ is a non-zero square modulo $r$, so its Legendre
symbol is $1$. For $c = 1$,
$\left(\frac{-2}{r}\right) = \left(\frac{-1}{r}\right)\left(\frac{2}{r}\right)$,
which is $+1$ exactly for $r \equiv 1, 3 \pmod 8$ and $-1$ for
$r \equiv 5, 7 \pmod 8$.
\end{proof}

Lemma~\ref{lem:qr} is the strongest prune we have, because it applies to a
single component in isolation: a pair $(p, 2a)$ can be rejected outright, with no
reference to the rest of the state. Taking $c = 1$ and $a = 1$:
\[
  \sigma(3^{2}) = 13,\quad \sigma(5^{2}) = 31,\quad
  \sigma(257^{2}) = 61 \cdot 1087,\quad \sigma(65537^{2}) = 37 \cdot 116085511,
\]
and $13 \equiv 5$, $31 \equiv 7$, $61 \equiv 5$, $37 \equiv 5 \pmod 8$. So none
of $3^{2}$, $5^{2}$, $257^{2}$, $65537^{2}$ can be the exact power of its prime
in a quasiperfect number, while $\sigma(17^{2}) = 307 \equiv 3$ is permitted.
Section~\ref{sec:fermat} draws the consequence for the hardest states.

Lemma~\ref{lem:qr} is not new, and we make no claim of priority. It was, however,
derived here independently: stated, proved, implemented, tested for soundness
against solutions with $c > 1$, and measured at a factor of $130$
(Section~\ref{sec:results}), before we found it in the literature as the
\emph{Legendre--Cattaneo filter} --- in exactly the form used here, including its
application to each component $\sigma(p^{2a})$, and formally verified in Lean~4 by
the UALBF project~\cite{ualbf}. It is \emph{not} the ``quadratic-residue sieve''
of~\cite{tty}, as we previously guessed it was: theirs tests whether a
discriminant is a quadratic residue to several small moduli, sieving candidate
primes before an exact perfect-square check, which is a condition on a different
object altogether (Section~\ref{sec:theirlemmas}).

That sequence is worth recording for two reasons beyond candour. First it is
checkable rather than merely asserted: the development history, the timestamped
source variants of Section~\ref{sec:variants} and the git history
together fix the order of events. Second, the convergence is evidence. An
independent derivation reaching a known classical result, in the same unusual
per-component form, corroborates both the derivation and the implementation
built on it --- the same sort of check as the two unrelated decompositions of
Section~\ref{sec:replication} agreeing on the exact state count, and rather
stronger, since the agreement here is with the literature rather than with
ourselves. What it does not do is make the lemma ours, and
Section~\ref{sec:history} records why we should have found it sooner.

\begin{remark}
In implementation the test is deliberately one-sided. Deciding it exactly needs
the full factorisation of $\sigma(p^{2a})$, which is infeasible for large $p$ and
$a$, so we factor only as far as is cheap and reject only when a violating prime
is actually exhibited. Failing to factor therefore costs pruning power and never
costs soundness. The cheap Jacobi test $\left(\frac{-2c}{\sigma}\right) = -1$ is
a useful first filter but is not equivalent: two violating primes can cancel, as
in $\sigma(257^{2}) = 61 \cdot 1087$ where both factors are bad yet the Jacobi
symbol of the product is $+1$.
\end{remark}

\subsection{A free corollary: the mod-8 exponent test}\label{sec:mod8}
Lemma~\ref{lem:qr} requires factoring $\sigma(p^{2a})$, which is why the
implementation caps the effort and so sometimes fails to reject a component it
should. A weaker form of it costs nothing at all.

\begin{corollary}\label{cor:mod8}
For $c = 1$, if $p^{2a}$ is the exact power of $p$ dividing $n$ then
\[
  \sigma(p^{2a}) \equiv 1 + a(p+1) \equiv 1 \text{ or } 3 \pmod 8 .
\]
\end{corollary}
\begin{proof}
The residues $\{1,3\}$ are closed under multiplication modulo $8$, since
$3 \cdot 3 = 9 \equiv 1$. By Lemma~\ref{lem:qr} every prime factor of
$\sigma(p^{2a})$ lies in that set, so the product does too. For the explicit
form, $p^{2} \equiv 1 \pmod 8$ for odd $p$, so of the $2a+1$ terms of
$\sigma(p^{2a})$ the $a+1$ even powers are each $\equiv 1$ and the $a$ odd powers
are each $\equiv p$, giving $\sigma(p^{2a}) \equiv (a+1) + ap = 1 + a(p+1)$.
\end{proof}

This is a test on $a$ and $p \bmod 8$ alone: no factoring, no large arithmetic.
The admissible exponents are $a \equiv 0, 1 \pmod 4$ for $p \equiv 1 \pmod 8$;
$a$ even for $p \equiv 3 \pmod 8$; $a \equiv 0, 3 \pmod 4$ for
$p \equiv 5 \pmod 8$; and no restriction for $p \equiv 7 \pmod 8$, where
$p + 1 \equiv 0$ and the congruence is vacuous. The identity was checked over
$2\,256$ pairs $(p, a)$.

Three of these are the exponent conditions of Cattaneo's 1951
theorem~\cite{cattaneo}, which we derived here before finding them there and which
we had in any case read only from a damaged extraction of that paper --- it
disagrees with the derivation for $p \equiv 5 \pmod 8$, and we follow the
derivation. Cattaneo's fourth condition, $a \equiv 1$ for $p \equiv 7 \pmod 8$,
does not follow from this argument even working modulo $16$, and we do not use
it.

That last point was checked rather than assumed, since modulo $16$ looked like
the obvious place to find the missing condition. It is not there, and the reason
is structural: the residues $\{1, 3, 9, 11\}$ modulo $16$ --- exactly those
congruent to $1$ or $3$ modulo $8$ --- are closed under multiplication, forming
a cyclic group generated by $3$. The argument above therefore reproduces itself
one power higher and yields nothing new per component, which we confirmed
numerically over more than two thousand pairs $(p, a)$: not one is rejected
modulo $16$ that survives modulo $8$. For $p \equiv 7 \pmod 8$ in particular,
$\sigma(p^{2a})$ is $1$ or $9$ modulo $16$ for every $a$, both admissible.

Working modulo $16$ does yield one thing, but a global one. An odd square is
$1$ or $9$ modulo $16$, so for $c = 1$
\[
  \sigma(n) = 2n + 1 \equiv 3 \pmod{16}
\]
exactly, for every candidate. Writing each component's level $\lambda(p, 2a)$
for the exponent with $3^{\lambda} \equiv \sigma(p^{2a}) \pmod{16}$, the levels
of a complete factorisation must sum to $1$ modulo $4$. This is a genuine
necessary condition and it is one we have not seen stated, but it is a single
equation over $\mathbb{Z}/4$ on the \emph{complete} component set, so it
constrains only the final component: a four-fold cut on leaves. As
Section~\ref{sec:wronglever} records, leaves are where this search spends
almost none of its time, so we verified the condition and did not implement it.

Being a consequence of Lemma~\ref{lem:qr} the corollary adds no reach in
principle, but it adds a great deal in practice, because it fires where the
factoring gives up: at $\omega = 6$ it cuts the state count from
$20\,211\,100$ to $15\,353\,292$ and the time from twenty-one minutes to eight.
It is guarded to $c = 1$: for other $c$ the admissible residues are those with
$\left(\frac{-2c}{r}\right) = 1$, which need not be closed under multiplication,
so the step from each factor to the product fails.

Lemma~\ref{lem:parity} is what makes refinement terminate. As the open
exponents of a state grow, $\lo$ increases to the limit on the left-hand side,
which is either below $2$ or above it. In the first case $\hi < 2$, and the
search may add a prime (or, when $r = 0$, prune the state). In the second,
$\lo$ eventually exceeds $2$, or $2 + c/n_{\min}$, and the state is pruned by
Lemma~\ref{lem:range}(4) or Lemma~\ref{lem:complete}.

\subsection{Order of operations}
At a state with $r \ge 1$: prune by Lemma~\ref{lem:range}(4) and~(3); if
$\hi \ge 2$, refine; if $r \ge 2$, some component is open and
$r\hi/(2-\hi) > 5000$ (so the next prime would range over a very long
interval), refine; if $r = 1$ and no component is open, close by
Lemma~\ref{lem:lastprime}; if $r = 1$ and the admissible range for $q$
(Lemma~\ref{lem:range}(2),(3)) is wider than a fixed threshold, refine;
otherwise add each admissible $q$, subject to Lemmas~\ref{lem:coprime}
and~\ref{lem:range}(1). Whenever a component is fixed at exponent $2a$, it must
pass Lemmas~\ref{lem:coprime} and~\ref{lem:qr}. At a state with $r = 0$: if
nothing is open, test
\eqref{eq:target} directly; otherwise prune by Lemma~\ref{lem:complete}, and
close by Lemma~\ref{lem:oneexp} if exactly one component is open, or refine the
open component with the least exponent if several are.

The threshold rule for $r = 1$ affects efficiency only. Its motivation: as the
open intervals tighten, the admissible range for the last prime shrinks to a
width close to $1$, so refining first avoids enumerating a wide range of last
primes with exponents still undetermined. Applying the same rule at $r \ge 2$
was tried and made the search roughly ten times slower.

\subsection{A worked example}\label{sec:worked}
The case $c = 1$, $k = 3$ shows the bounds doing their work, and is small enough
to check by hand. Write $f(p) = p/(p-1)$.

At the root, $r = 3$ and Lemma~\ref{lem:range}(3) requires
$f(p_1)f(p_2)f(p_3) > 2$ over the three smallest admissible primes. For
$p_1 = 3$ this is $f(3)f(5)f(7) = \tfrac32\cdot\tfrac54\cdot\tfrac76
= \tfrac{105}{48} = 2.1875 > 2$, so $p_1 = 3$ survives; for $p_1 = 5$ it is
$f(5)f(7)f(11) = \tfrac{385}{240} \approx 1.604 < 2$, so every branch with
$p_1 \ge 5$ is dead at once. Hence $3 \mid n$.

Enter $3$ as open, so $\lo = h(3^2) = \tfrac{13}{9}$, $\hi = f(3) = \tfrac32$,
$n_{\min} = 9$, and $r = 2$. Since $\hi < 2$ no refinement is forced. For the
next prime $q$, Lemma~\ref{lem:range}(3) now requires
$\hi\, f(q) f(q') > 2$ with $q'$ the prime after $q$: for $q = 5$,
$\tfrac32\cdot\tfrac54\cdot\tfrac76 = 2.1875 > 2$, but for $q = 7$,
$\tfrac32\cdot\tfrac76\cdot\tfrac{11}{10} = 1.925 < 2$. So $q = 5$ is forced and
$15 \mid n$.

Now $\lo = \tfrac{13}{9}\cdot\tfrac{31}{25} = \tfrac{403}{225}$,
$\hi = \tfrac32\cdot\tfrac54 = \tfrac{15}{8}$, $n_{\min} = 225$, $r = 1$. The
last prime is confined from below by Lemma~\ref{lem:range}(2),
\[
  q > \frac{\lo}{2 - \lo + 1/n_{\min}}
    = \frac{403/225}{(450 - 403 + 1)/225} = \frac{403}{48} \approx 8.40,
\]
and from above by Lemma~\ref{lem:range}(3), which gives
$f(q) > 2/\hi = \tfrac{16}{15}$, that is $q < 16$. Only $q = 11$ and $q = 13$
remain, and both pass the sharper test of Lemma~\ref{lem:range}(1) and the
coprimality test of Lemma~\ref{lem:coprime} (neither divides
$\sigma(3^2) = 13$ or $\sigma(5^2) = 31$). Each of the two surviving states is
then refined until Lemma~\ref{lem:parity} forces $\hi$ below $2$ or $\lo$ above
$2 + 1/n_{\min}$, and both die. The whole of $\omega = 3$ costs eight states.

The pattern generalises: the upper bound of Lemma~\ref{lem:range}(3) makes the
smallest primes almost forced, while the admissible window for the last prime
is what grows explosively when the fixed part's abundancy approaches $2$. That
is the mechanism behind the Fermat-prime subtree of Issue~7.

% ---------------------------------------------------------------------------
\section{Verification}\label{sec:verification}

The implementation is about 300 lines of Python using exact rational and
integer arithmetic throughout. Primality of candidate primes is decided by
Miller--Rabin with the first twenty prime bases. This can only err by calling a
composite prime, which adds a spurious branch and cannot remove a genuine one.

\subsection{Brute force}
For every odd $k \le 30001$ we computed $c = \sigma(k^2) - 2k^2$. For each pair
$(c, \omega)$ with $0 < c \le 2 \times 10^5$ found this way (thirteen pairs,
with $\omega \in \{3, 4\}$), the search was run with that $c$ and $\omega$.
In every case the solutions it returned that were at most $30001^2$ coincided
exactly with the brute-force list, and no state was left unresolved.

\subsection{Planted solutions}
Brute force cannot reach large exponents or large last primes, which is exactly
where the closing lemmas operate. We therefore built random odd squares
$n = \prod p_i^{2a_i}$ with two to five primes, exponents $2a_i$ up to $24$
and a largest prime up to $2000$. For each $n$ with $c = \sigma(n) - 2n > 0$,
the closing routines were asked to recover the missing piece of $n$:
\begin{itemize}
\item \emph{Last prime}: all components but the largest prime fixed, whenever
  that remainder satisfies $\sigma(m) < 2m$ (Lemma~\ref{lem:lastprime}).
\item \emph{One exponent}: all primes known and one exponent, chosen at
  random, open from $2$ (Lemma~\ref{lem:oneexp}).
\end{itemize}
Across four random seeds, $8\,000$ instances of each type were recovered with
no failures. The planted values of $c$ are far larger than $1$, which
exercises parts of the last-prime root finding that never arise for $c = 1$.

\subsection{Internal consistency}
The parallel driver, which cuts the tree at a fixed depth into independent
subtrees, reproduces the serial node count exactly for $\omega = 5$
($193\,559$ nodes). The tool that splits a single subtree further
(Issue~7) was checked the same way, against an already completed subtree whose
node count it reproduced exactly ($200\,601$).

\subsection{Soundness of the quadratic-residue prune}\label{sec:qrtest}
Lemma~\ref{lem:qr} is the only prune that rejects a component without reference
to the rest of the state, so an error in it would silently remove solutions
rather than produce a visible unresolved state. It was therefore tested three
ways.

First, directly: the lemma predicts that $13$, $31$, $61$ and $37$ never divide
$2m^{2}+1$, while $11$ and $307$ (both $\equiv 3 \bmod 8$) may. Over
$m \le 2 \times 10^{5}$ the former divide it not once --- indeed $0$ is not among
the residues of $2m^{2}+1$ modulo $13$ --- and over $m \le 2\times 10^{4}$ the
latter divide it $3\,636$ and $130$ times respectively.

Second, against real solutions. The lemma generalises to arbitrary $c$, which
matters because solutions exist for $c > 1$ and can be found by brute force. Of
the $37$ genuine solutions of $\sigma(n) = 2n + c$ with $c \le 10^{6}$ and
$m \le 8000$, and all $131$ of their components, not one is rejected. Any
rejection would have been a counterexample.

Third, inside the search. With the prune enabled the brute-force comparison of
Section~\ref{sec:verification} still recovers every known solution for all
thirteen $(c, \omega)$ pairs, with no mismatches and nothing unresolved, and the
planted-leaf suites still recover all their instances. Node counts drop while
results do not change, which is the signature wanted of a sound prune.

\subsection{Testing the faulty versions}
The defects described in Section~\ref{sec:history} are claims about particular
versions of the program, so they are tested as such. Each recovered faulty
variant (Section~\ref{sec:variants}) is run against the defect attributed to it,
alongside its replacement on the same input; the results appear with the issues
themselves. This is not redundant with the suites above, which test only the
current version: it is what stops the narrative of Section~\ref{sec:history} from
being unfalsifiable. In one case it showed our first test was measuring nothing
(Issue~8).

% ---------------------------------------------------------------------------
\section{Results}\label{sec:results}

\begin{table}[h]
\centering
\begin{tabular}{cccc}
\toprule
$\omega$ & nodes & solutions & unresolved states \\
\midrule
1 & 1 & 0 & 0 \\
2 & 1 & 0 & 0 \\
3 & 8 & 0 & 0 \\
4 & 542 & 0 & 0 \\
5 & 193\,559 & 0 & 0 \\
6 & 1\,998\,951\,360 & 0 & 0 \\
7 & partial only; see Section~\ref{sec:omega7cost} & & \\
\bottomrule
\end{tabular}
\caption{Search for $\sigma(n) = 2n + 1$ with $n$ an odd square, by number of
distinct prime factors, \emph{without} Lemma~\ref{lem:qr}. A run establishes its
row only if the unresolved column is zero.}
\end{table}

\begin{table}[h]
\centering
\begin{tabular}{crrr}
\toprule
$\omega$ & states without Lemma~\ref{lem:qr} & with it & reduction \\
\midrule
3 & 8 & 7 & $1.1\times$ \\
4 & 542 & 208 & $2.6\times$ \\
5 & 193\,559 & 15\,176 & $12.8\times$ \\
6 & 1\,998\,951\,360 & 15\,353\,292 & $130\times$ \\
\bottomrule
\end{tabular}
\caption{Effect of Lemma~\ref{lem:qr} together with its free corollary, the
mod-8 exponent test of Section~\ref{sec:mod8}. Every row still reports no
solution and no unresolved state. The reduction grows with $\omega$, as one would
expect of a test applied independently to each component.}
\end{table}

Both runs of each row agree: no solutions, nothing unresolved. So $\omega = 6$
has now been settled three times over by this implementation --- once as $358$
parallel subtrees, once with its hardest subtree decomposed a second way
(Section~\ref{sec:replication}), and once in a single process with the prunes
enabled. The last took eight minutes on one core against $17.4$ hours on
fourteen, a reduction of about $288$ in total processor time. Notably its exact
leaf closures fell from $335$ to zero: with the prune in place, no branch
survives long enough to need one.

\begin{theorem}\label{thm:main}
No quasiperfect number has six or fewer distinct prime factors.
\end{theorem}

This is not new: it is one of the two bounds of Hagis and
Cohen~\cite{hagiscohen}, and its value here lies only in the independence of the
method and the code. It is worth being exact about how much of that paper is
reproduced, which is less than the phrase ``the 1982 bound'' suggests.
Theorem~\ref{thm:main} corresponds to their prime-factor bound alone. Their size
bound $n > 10^{35}$ rests on different techniques and is untouched by anything
here; since raised to $10^{45}$ by Alekseyev~\cite{alekseyev}, it is the natural
next target and is discussed in Section~\ref{sec:further}. We have also not
examined whatever subsidiary results their paper contains beyond these two, so
no claim is made about them. And the search assumes throughout that a
quasiperfect number is an odd square, which is Cattaneo's
theorem~\cite{cattaneo}, cited and not reproved
(Lemma~\ref{lem:oddsquare} and the remark following it). Our result is therefore
conditional on that one prior result holding.

\subsection{The shape of the $\omega = 6$ computation}
The tree was cut at four primes into $358$ independent subtrees. The work is
extraordinarily lopsided: $208$ of them finish in under ten seconds, the
slowest of the ordinary ones takes $2.6$ hours, and a single subtree ---
the one beginning $3, 5, 17, 257$ of Issue~7 --- accounts for
$1.405 \times 10^{9}$ states, just over $70\%$ of the whole computation.
Together the run took about $17.4$ hours of wall-clock time and, counting each
task once, some $38$ core-hours. Only $335$ exact leaf closures were needed in
the whole of $\omega = 6$: almost everything dies to the bounds of
Lemma~\ref{lem:range} long before a leaf is reached.

Coverage was audited separately from the run, by re-deriving the list of
subtrees from the splitter and checking that every one has a record, that no
record belongs to a subtree outside the list, and that no record reports a
solution or an unresolved state. That audit passes, which is what turns
``nothing was found'' into Theorem~\ref{thm:main} via
Proposition~\ref{prop:complete}.

\subsection{An unplanned replication}\label{sec:replication}
The concurrency of Issue~9 produced a stronger check than we designed. The
Fermat subtree was searched twice to completion by different routes: once as a
single process walking it whole, and once as $12\,195$ separate processes, each
handling one candidate fifth prime, with the enumeration of that prime lifted
into the splitter. The two runs report exactly the same number of search
states, $1\,404\,698\,494$. Since the decomposition is textually unrelated to
the traversal --- different process structure, different order, different entry
points --- an agreement to the last digit is good evidence that the tree is
well defined and that the implementation walks it deterministically.

The two routes are not equally efficient, and the difference measures Issue~8
directly. The split consumed $218\,103$ core-seconds against the single
process's $62\,605$: a factor $3.5$ more work for the same states, because each
of the $12\,195$ pieces re-refines the exponents of $3, 5, 17, 257$ from
scratch. In exchange it finished in $16\,779$ seconds of wall-clock time
against $62\,605$, a speedup of $3.7$ on thirteen cores ($29\%$ efficiency).
So the split bought a useful reduction in wall-clock time at a large cost in
total work, and the refinement rule of Issue~8 addresses exactly the waste it
introduced.

% ---------------------------------------------------------------------------
\section{Conditional bounds}\label{sec:conditional}

The results of this section may be new. They are also the cheapest computations
in the paper, for a reason that is worth stating plainly: barring small primes is
exactly the hypothesis that makes this method work well, and the literature
appears not to have searched those cases beyond the elementary count.

\subsection{Why the restricted cases are easy}
Everything expensive in Sections~\ref{sec:fermat} and~\ref{sec:omega7cost} came
from configurations in which $\prod p/(p-1)$ sits within $10^{-5}$ of $2$, which
requires the small primes. Remove them and the arithmetic changes character.
Reading Lemma~\ref{lem:range}(3) as a bound on $\omega$ by restricting the
available primes:

\begin{table}[h]
\centering
\begin{tabular}{lrl}
\toprule
primes barred & $\omega \ge$ & best published \\
\midrule
none      & $3$  & $7$~\cite{hagiscohen}, $8$ claimed~\cite{tty} \\
$5$       & $4$  & none found; the unconditional bound applies \\
$3$       & $7$  & none found; the unconditional bound applies \\
$3$ and $5$ & $15$ & $15$~\cite{prasad}, refining $10$ (Abbott et al., 1973) \\
\bottomrule
\end{tabular}
\caption{Primes needed for $\prod p/(p-1) > 2$ once small primes are barred,
against the best bounds we can find in the literature. Note that
$\omega \ge 15$ for $\gcd(n,15) = 1$ is exactly the elementary count.}
\end{table}

The published record for the restricted cases, taken from the introduction
of~\cite{prasad} rather than from any summary of it, is: $\omega \ge 5$ (Abbott,
Aull, Brown and Suryanarayana, Theorem 2), improved to $6$ by Kishore and to $7$
by Cohen and Hagis; and conditionally $\omega \ge 10$ if $\gcd(N,15) = 1$ (the
same authors, Theorem 5), refined to $15$ by~\cite{prasad}, who also give an
alternative proof of the unconditional bound when $\gcd(N,15) = 5$. We can find
no bound stated for $3 \nmid n$ with $5$ permitted, so for that hypothesis the
unconditional bound is what stands.

The last row is the opening. Unconditionally, search has bought four or five
beyond the elementary bound; in the coprime-to-$15$ case the published bound
\emph{is} the elementary bound, so no search appears to have been applied there
at all.

\subsection{Results}
Running the search of this paper with those primes barred gives the following.
Each line means the case was eliminated with no solution found and no state left
unresolved.

\begin{table}[h]
\centering
\begin{tabular}{llrrr}
\toprule
hypothesis & $\omega$ eliminated & states & new bound & previous \\
\midrule
$5 \nmid n$          & $4$  & $24$ & & \\
                     & $5$  & $703$ & & \\
                     & $6$  & $41\,267$ & & \\
                     & $7$  & $19\,899\,890$ & $\omega \ge 8$ & $7$ \\
\midrule
$3 \nmid n$          & $7$  & $15$ & & \\
                     & $8$  & $1\,871$ & & \\
                     & $9$  & $540\,250$ & $\omega \ge 10$ & $7$ \\
\midrule
$\gcd(n,15) = 1$     & $15$ & $2\,159$ & & \\
                     & $16$ & $788\,607$ & & \\
                     & $17$ & $1.51 \times 10^{10}$\textsuperscript{$*$} & $\omega \ge 18$ & $15$~\cite{prasad} \\
\bottomrule
\end{tabular}
\caption{Conditional bounds. Values of $\omega$ below those listed are excluded
by the elementary count of the previous table, in one state each. Cost grows by
roughly a factor of $300$ per step, so one or two further steps in each case is
the practical limit. ${}^{*}$\,$12\,916\,647\,713$ states of the exact search
plus $2\,199\,616\,513$ of the compiled search that closed its one unresolved
node (Section~\ref{sec:omega17}).}
\end{table}

These improve the bound for $\gcd(n,15) = 1$ from $15$ to $18$, and give $10$ for
$3 \nmid n$ and $8$ for $5 \nmid n$, where only the unconditional bound of $7$
previously applied. The last of those is the expensive one, at twenty million
states against a few hundred thousand for the others, because permitting $3$
readmits the configurations that drive the abundancy towards $2$.

\subsection{The last step for $\gcd(n,15) = 1$: closing one node}\label{sec:omega17}
The search for $\omega = 17$ with $3$ and $5$ barred was split into $507$
subtrees and run between 2 and 11 October 2026. It took $12\,916\,647\,713$
states and found no solution, the last four subtrees each taking between $38$
and $152$ hours on one core. It did not finish cleanly. One state in the last
subtree was left unresolved: fifteen components, nine of them with open
exponents, and two primes $q < q'$ still to choose. Enumeration of $q$ stopped
at the limit of $10^{7}$ candidates, at $q = 571\,287\,683$, while the
admissible range runs to $1\,137\,878\,224$, leaving $27\,577\,412$ primes.

\paragraph{Why the remainder is dear.} Write $\eta = 2 - \mathrm{hi}$ for the
node; here $\eta = 3.515 \times 10^{-9}$. Below $q = 2/\eta \approx 5.69 \times
10^{8}$ each $q$ forces some open exponents to stay small, and the first ten
million $q$ cost about $450$ states each. Above it every open exponent may grow
without bound for every $q$, and the exact search spends about $8\,000$ states
and two seconds per $q$ re-refining the same nine exponents: roughly
$15\,000$ core-hours for the remainder. The run had stopped almost exactly
where its cost changes character.

\paragraph{Issue: a dead end.} Since neighbouring $q$ give almost the same
tree, we first carried $q$ as an interval of primes, with bounds valid for
every prime in it, and split the interval only when its own spread dominated
(\texttt{qblock.py}). It agreed with the exact search on the test suites and
was slower in practice: the surviving exponent boxes are numerous at every
scale, so the shared upper levels saved little and the late splits repeated
the work. The cost lies in the number of exponent patterns per $q$, not in
the arithmetic of each.

\paragraph{The fix: the same search, compiled, in logarithms.} An open
component $p$ with exponent $2a$ contributes $\log\frac{p}{p-1} - d_p(a)$ to
the log-abundancy, where $d_p(a) = -\log(1 - p^{-(2a+1)})$, so the equation
becomes
\[
  \log f(q) + \log f(q') \;=\; K + \sum_p d_p(a_p) + \log\Bigl(1 + \frac{1}{2N}\Bigr),
  \qquad K = \log\frac{2}{\mathrm{hi}},
\]
where $\log f(x) = \log\bigl(\sigma(x^{2b})/x^{2b}\bigr)$ lies in
$[\log(1 + 1/x + 1/x^{2}),\ -\log(1 - 1/x))$ for every $b \ge 1$. For each $q$
the program (\texttt{qfast.py}, compiled with Numba) keeps a box over the open
exponents, each either fixed or open from a floor. It converts the box into
the range of $q'$ it permits, and splits the open component with the widest
contribution until that range is short. It then tries each prime $q'$ and
splits further until the box is narrower than the target the pair leaves. The
residue test on fixing an exponent is kept, precomputed per $(p, a)$ with the
exact search's own function. The other prunes are dropped, which can only
enlarge the tree.

\paragraph{Issue: precision.} The terms are about $1.8 \times 10^{-9}$, and a
first version in double precision needed a slack of $2 \times 10^{-23}$ on every
comparison. It left about $0.56$ survivors per $q$ for the exact code at
$75$\,ms each, some $90$ core-hours in all. Moving the sums to double-double
arithmetic brought the rounding error to about $10^{-41}$; $K$ and every
$d_p(a)$ are summed from their series as exact rationals and then split into
high and low parts. Checked against exact rationals, the computed bounds on
$\log f(x)$ are within $1.3 \times 10^{-41}$. The slack is $10^{-35}$, so a
comparison can only keep too much, never discard a solution. Survivors fell
to $22$ in the whole range.

\paragraph{Result.} All $27\,577\,412$ primes $q$ in
$(571\,287\,683,\ 1\,137\,878\,224]$ took $2\,199\,616\,513$ states in about
fifteen minutes on one core. The $22$ survivors were handed, as boxes, to the
original exact search, which settled them in $1\,418$ states with no solution
and nothing unresolved. With the node closed, $\omega = 17$ is eliminated when
$\gcd(n,15) = 1$, so $\omega(n) \ge 18$.

\paragraph{Validation, and its limit.} The compiled search was checked by
planting. For $300$ random choices of $q$, of $q'$ up to $50q$, and of a fixed
exponent vector, $K$ was set so that the planted point satisfies the equation
exactly; every time, a returned box contained it. That tests the splitting and
the bounds. It does not test the computation of $K$ for the real node, which
rests on the rational series above. As a second, independent check, the
original exact code was run on $4\,000$ primes sampled at random from the same
range: $13\,172\,539$ states, no solution and nothing unresolved
(\texttt{spot\_qrange.py}). The
compiled search is new code, checked only by planting and sampling, and an
independent re-run of this node is the check we would most like to see.

\paragraph{Issue: two slips while running it.} Our first estimate for the
remainder was ``hours'' by the discriminant criterion. That was wrong: the
criterion needs fixed exponents, and nine were open. And the first full launch
spent eighty minutes doing nothing. The end of the $q$-range had been written
as $1/(1-s)$ instead of $s/(s-1)$, with $s = \sqrt{2/\mathrm{hi}}$; the result
is negative, and a loop meant to adjust the end by a few units counted downwards
indefinitely. The block-search version had the formula right, and it was
retyped rather than reused.


\subsection{A consequence for the unconditional problem}\label{sec:forced15}
The two one-prime cases combine into a statement with no hypothesis attached.

\begin{proposition}\label{prop:forced15}
Subject to the implementation (Section~\ref{sec:conditional-caveats}), any
quasiperfect number with exactly seven distinct prime factors is divisible by
$15$.
\end{proposition}
\begin{proof}
If $3 \nmid n$ then $\omega(n) \ge 10$, so $\omega(n) = 7$ forces $3 \mid n$. If
$5 \nmid n$ then $\omega(n) \ge 8$, so $\omega(n) = 7$ forces $5 \mid n$.
\end{proof}

Since $\omega(n) \ge 7$ unconditionally~\cite{hagiscohen}, $\omega = 7$ is the
smallest case not already excluded, and this says the whole of it lies under
$15 \mid n$.

\medskip\noindent\textbf{But note the vacuity risk.} If the claim
$\omega(n) \ge 8$ of~\cite{tty} is correct then no quasiperfect number has seven
prime factors at all, and Proposition~\ref{prop:forced15} is true but empty. We
therefore do not present it as a new theorem. Its value is as \emph{partial
independent verification} of that unverified claim: eliminating $\omega = 7$
entirely is what~\cite{tty} does, and we have eliminated the part of it with
$3 \nmid n$ or $5 \nmid n$, by a method of our own, at a cost of twenty million
states. What remains unchecked is $\omega = 7$ with $15 \mid n$. We have also not
read~\cite{hagiscohen} in full and cannot say whether a result of this shape is
already there.

\medskip\noindent There is an irony in the arithmetic. The cases
Proposition~\ref{prop:forced15} leaves standing --- those with $3$ and $5$ both
dividing $n$ --- are exactly the configurations of
Section~\ref{sec:fermat} that make $\prod p/(p-1)$ approach $2$ and that consumed
essentially the whole of our $\omega \le 6$ computation. So the structural result
and the intractability of unconditional $\omega = 7$ are two descriptions of one
fact: every configuration that is cheap to search is excluded, and the residue is
precisely what is dear.

\subsection{Why this was not done before}\label{sec:whynot}
The obvious objection to these results is that they were cheap, so somebody would
have had them. Three things seem to account for it, and the first two are visible
in the literature rather than inferred.

The restricted cases have been approached \emph{analytically}, not
computationally. Abbott et al.\ derived $\omega \ge 10$ for $\gcd(N,15) = 1$ in
1973 and~\cite{prasad} refined it to $15$ in 2017 by proving a structural theorem
about what they call special factors. Neither is a search. Exhaustive search of
the conditional cases was not available in 1973, and by the time it was, the
attention had moved to the unconditional problem, where it is genuinely hard.

Second, the tools needed are old. Cattaneo's 1951 theorem already constrains the
exponents by the residue class of each prime modulo $8$, and the
Legendre--Cattaneo filter of Lemma~\ref{lem:qr} dates from the same literature.
Nothing used here is modern except the arithmetic throughput and the willingness
to spend it.

Third, and the reason the cost is so lopsided, the restricted hypotheses remove
exactly what makes the unconditional problem expensive. Everything difficult in
Sections~\ref{sec:fermat} and~\ref{sec:omega7cost} needs $\prod p/(p-1)$ within
about $10^{-5}$ of $2$, which needs the small primes. Barring $3$ and $5$ makes
the smallest admissible prime $7$ and forces $7, 11, 13, 17, 19, 23, 29, 31$ in
order with no branching at all. The conditional cases are not hard problems that
we solved; they are easy problems that nobody had run.

We also note, since we are not using it, that Cattaneo's exponent conditions
would prune the search further. Our results are unaffected --- omitting a
necessary condition searches a superset and so cannot create a false
elimination --- but they could be considerably faster, and that is the obvious
way to push these bounds another step. The depth profile of the latter is
instructive: of the $2\,159$ states at $\omega = 15$, there is exactly one at each
of the first eight levels --- the primes $7, 11, 13, 17, 19, 23, 29, 31$ are
forced in order before any branching is possible --- and $2\,070$ of them sit at
the final level.

\subsection{What these rest on}\label{sec:conditional-caveats}
More than usually, the reader is asked to take the implementation on trust, so we
state the support precisely. The search is exhaustive by
Proposition~\ref{prop:complete} and terminating by
Proposition~\ref{prop:terminate}; it reproduces Theorem~\ref{thm:main}, a result
with a known answer, three times over; and it passes the suites of
Section~\ref{sec:verification}. The prime-barring mechanism was checked in both
directions: with $3$ and $5$ barred the search correctly returns nothing for
$c = 921$, whose only solution is $3^{2}5^{2}7^{2}$, and it does not over-bar ---
each of $13$ known solutions is still recovered when a prime that does not divide
it is barred. Nothing here has been verified by anyone else.

We also note that the UALBF project~\cite{ualbf} has open work on a bound of the
form $\omega \ge 11$ in a case involving $5$, so this ground is being worked
concurrently and these bounds may be matched or exceeded elsewhere.

% ---------------------------------------------------------------------------
\section{The Fermat-prime obstruction}\label{sec:fermat}

The subtree that dominated $\omega = 6$ was not an accident of our search order.
This section identifies the structure behind it, which appears to be the same
obstruction that~\cite{tty} describes as a ``family of deep leaves on which pure
enumeration cannot terminate''. We have read only the abstract of that paper, so
we cannot say whether the characterisation below is how its authors describe the
difficulty; we state it because it is what our own computation forced on us, and
because it predicts precisely where $\omega = 7$ becomes hard.

Write $F_i = 2^{2^i} + 1$, so $F_0, \dots, F_4 = 3, 5, 17, 257, 65537$ are the
known Fermat primes.

\begin{lemma}[The greedy chain is the Fermat chain]\label{lem:fermat}
For $k \ge 0$,
\[
  P_k := \prod_{i=0}^{k} \frac{F_i}{F_i - 1} = 2 - 2^{\,1 - 2^{k+1}} .
\]
Moreover, if $q$ is a prime with $P_k \cdot q/(q-1) < 2$, then
$q > 2^{2^{k+1}} - 1$; so the least such candidate is $F_{k+1}$ itself, which is
available only when $F_{k+1}$ is prime.
\end{lemma}
\begin{proof}
Since $(x-1)\prod_{i=0}^{k}(x^{2^i}+1) = x^{2^{k+1}} - 1$, putting $x = 2$ gives
$\prod_{i=0}^{k} F_i = 2^{2^{k+1}} - 1$, while
$\prod_{i=0}^{k}(F_i - 1) = \prod_{i=0}^{k} 2^{2^i} = 2^{2^{k+1}-1}$. Dividing
gives $P_k = (2^{2^{k+1}}-1)/2^{2^{k+1}-1} = 2 - 2^{1-2^{k+1}}$. The condition
$P_k\, q/(q-1) < 2$ rearranges to $q > P_k/(2 - P_k) = 2^{2^{k+1}} - 1$.
\end{proof}

So the sequence obtained by repeatedly taking the least odd prime that keeps
$\prod p/(p-1)$ below $2$ is exactly $3, 5, 17, 257, 65537$, with the thresholds
$1, 3, 15, 255, 65535$, each one less than a power of two. The chain terminates
there, not by choice but because $F_5 = 2^{32}+1 = 641 \cdot 6700417$ is
composite. These five primes therefore give the abundancy that approaches $2$
from below more closely than any other set of five odd primes, reaching
$2 - 2^{-31}$, and by Lemma~\ref{lem:range}(3) an abundancy that close to $2$ is
exactly what makes the admissible window for the remaining primes enormous.

\subsection{Why \texorpdfstring{$\omega = 6$}{omega = 6} was tractable and
\texorpdfstring{$\omega = 7$}{omega = 7} may not be}
Fixing each Fermat prime at exponent $2$ and applying
Lemma~\ref{lem:range}(2),(3) gives the following windows for the next prime.

\begin{table}[h]
\centering
\begin{tabular}{llll}
\toprule
$\omega$ & primes fixed & $r$ & window for the next prime \\
\midrule
6 & $3,5,17,257$          & 2 & $(21,\ 1.31 \times 10^{5})$ \\
7 & $3,5,17,257$          & 3 & $(21,\ 1.97 \times 10^{5})$ \\
6 & $3,5,17,257,65537$    & 1 & $(21,\ 4.29 \times 10^{9})$ \\
7 & $3,5,17,257,65537$    & 2 & $(21,\ 8.59 \times 10^{9})$ \\
\bottomrule
\end{tabular}
\end{table}

Two things follow. First, the four-prime states barely worsen from $\omega = 6$
to $\omega = 7$: the window grows only by half. They are not the difficulty.
Second, the five-prime states carry a window of order $2^{32}$ in both cases ---
but at $\omega = 6$ only one prime remains, and one unknown prime is closed in
\emph{closed form} by Lemma~\ref{lem:lastprime} with no enumeration at all. That
is why $\omega = 6$ finished. At $\omega = 7$ the same state has two primes
outstanding, and we have no closed form for two. The step from $r = 1$ to
$r = 2$ at a window of $2^{32}$ is the entire cliff, and it is consistent with
the scan bound of $8 \times 10^{8}$ for an ``intermediate prime'' quoted
in~\cite{tty} (whether that figure is a prime value or a count of primes is not
clear from the abstract; our window holds about $3.8 \times 10^{8}$ primes above
$65537$).

\subsection{The hard set is small}
The difficulty is confined much more tightly than ``states involving Fermat
primes''. Raising a single exponent from $2$ to $4$ in the five-prime,
$r = 2$ state changes the window as follows.

\begin{table}[h]
\centering
\begin{tabular}{lll}
\toprule
exponent raised & window & shrinks by \\
\midrule
$3^{2} \to 3^{4}$         & $4.86 \times 10^{2}$ & $1.8 \times 10^{7}$ \\
$5^{2} \to 5^{4}$         & $6.25 \times 10^{3}$ & $1.4 \times 10^{6}$ \\
$17^{2} \to 17^{4}$       & $2.84 \times 10^{6}$ & $3.0 \times 10^{3}$ \\
$257^{2} \to 257^{4}$     & $8.56 \times 10^{9}$ & $1.0$ \\
$65537^{2} \to 65537^{4}$ & $8.59 \times 10^{9}$ & $1.0$ \\
\bottomrule
\end{tabular}
\end{table}

So the hard states are those with $3^{2} \cdot 5^{2} \cdot 17^{2}$ dividing $n$
exactly, times $257$ and $65537$ to any power: a handful of states rather than a
family. Any increase in the exponent of $3$, $5$ or $17$ collapses the window by
three to seven orders of magnitude, while the exponents of the two large Fermat
primes are irrelevant to it.

\subsection{What Lemma~\ref{lem:qr} does and does not remove}
\label{sec:fermat-alive}
Lemma~\ref{lem:qr} forbids four of the five Fermat primes from appearing at
exponent exactly $2$, since $\sigma(3^{2}) = 13 \equiv 5$,
$\sigma(5^{2}) = 31 \equiv 7$, $\sigma(257^{2}) = 61 \cdot 1087$ and
$\sigma(65537^{2}) = 37 \cdot 116085511$ each contain a prime $\equiv 5$ or
$7 \pmod 8$; only $\sigma(17^{2}) = 307 \equiv 3$ is permitted. So the state with
all five at exponent $2$ --- the one that consumed $70\%$ of the $\omega = 6$
computation --- is excluded outright, four times over.

It is tempting to conclude that the obstruction is gone, and an earlier draft of
this paper did so. Measurement refutes it. The lemma bars exponent $2$ but
permits $3^{4}$, $3^{12}$, $5^{6}$, $5^{16}$ and so on, and those bring
$\prod p/(p-1)$ nearly as close to $2$ as exponent $2$ did. Of the ten dearest
depth-five subtrees in our $\omega = 7$ sample, eight have the shape
\[
  3^{a} \cdot 5^{b} \cdot 17^{c} \cdot 257^{d} \cdot q^{2},
  \qquad (a,b,c,d) \approx (10, 8, 8, 8), \quad q \sim 10^{5},
\]
with $2 - \lo$ of order $2 \times 10^{-5}$. The dominant cost at $\omega = 7$
therefore comes from the same five primes as at $\omega = 6$, at different
exponents.

The error in the optimistic reading is worth naming, because the table of the
previous subsection invites it. That table varies \emph{one} exponent while
holding the others at $2$, and in that setting raising the exponent of $3$ does
collapse the window by seven orders of magnitude. It does not follow that raising
several together is equally harmless, and it is not: the exponents above sit at
$(10,8,8,8)$ and the window stays wide. A one-at-a-time sensitivity analysis is
not a guide to joint variation.

What the prune achieved is a reduction in scale rather than an elimination. The
admissible window for the next prime in these states falls from
$4.3 \times 10^{9}$ to around $3 \times 10^{5}$, a factor of roughly $10^{4}$,
which accounts for most of the $130$-fold saving of
Section~\ref{sec:results} and for the $\omega = 7$ tree being $8.7$ times
narrower at depth five ($80\,458$ subtrees against $699\,207$). The obstruction is
ten thousand times smaller and in the same place, so progress at $\omega = 7$
requires attacking it directly --- which is what
Section~\ref{sec:theirlemmas} turns out to supply.

Lemma~\ref{lem:coprime} contributes a little more: with each Fermat prime
squared, none of $13$, $31$, $37$, $61$, $307$, $1087$, $116085511$ may divide
$n$. The exclusion of $37$ is the useful one, being small enough to have been an
attractive candidate.

\subsection{Exponents surviving the prune}
Exhaustive factorisation of $\sigma(p^{2a})$ for small $p$ and $c = 1$ leaves
remarkably few admissible exponents. This is the other reason the prune is
effective: it attacks the exponent dimension, which is the second of the two
obstacles that~\cite{tty} reports.

\begin{table}[h]
\centering
\begin{tabular}{ll}
\toprule
$p$ & exponents $2a \le 20$ surviving Lemma~\ref{lem:qr} \\
\midrule
$3$  & $4,\ 12$ \\
$5$  & $6,\ 16$ \\
$7$  & $2,\ 4,\ 10,\ 12,\ 16,\ 18$ \\
$11$ & $16$ \\
$13$ & $6$ \\
$17$ & $2,\ 8,\ 10$ \\
$19$ & $16$ \\
$23$ & $4,\ 10,\ 12$ \\
$29$ & none \\
$31$ & $2,\ 6,\ 16,\ 18$ \\
$37$ & none \\
$41$ & $2$ \\
$43$ & $4,\ 12$ \\
$47$ & $6,\ 10$ \\
$257$ & $16$ \\
\bottomrule
\end{tabular}
\caption{Admissible exponents for small primes. A prime admitting no exponent at
all cannot divide a quasiperfect number; $29$ and $37$ admit none up to $20$,
although ruling them out entirely would require treating every exponent, which
we have not done.}
\end{table}

% ---------------------------------------------------------------------------
\section{The cost of \texorpdfstring{$\omega = 7$}{omega = 7}, and why we cannot
state it}\label{sec:omega7cost}

\emph{This section records the costing as it stood before the compiled search of
Section~\ref{sec:omega7}, which completed $\omega = 7$ in about $61$ core-hours.
It is kept because the failure to price the search was itself informative.}

With Lemma~\ref{lem:qr} in hand, $\omega = 7$ no longer looked obviously out of
reach, so we tried to price it. We failed, and the manner of the failure is
worth recording, because it is a property of the problem rather than of our
effort.

\subsection{A uniform sample, and why it was not enough}
Sampling $400$ of the $80\,458$ depth-five subtrees under a cap of
$2 \times 10^{6}$ states each gave $2.93 \times 10^{8}$ states in $24\,846$
core-seconds, at $11\,795$ states per core-second, scaling to at least
$5.9 \times 10^{10}$ states and $1\,390$ core-hours --- about four days on
fourteen cores. That figure is worthless as an estimate: $122$ of the $400$
subtrees reached the cap and supplied $83\%$ of the sampled states while
truncated, so the quantity being measured was precisely the one the cap hid.

\subsection{A stratified sample}
We therefore stratified. A probe pass over $240$ subtrees at the same cap found
$71.2\%$ light, with an exact mean of $201\,739$ states, and $28.7\%$ heavy. A
second pass ran $28$ of the heavy ones at a cap fifty times larger,
$10^{8}$ states. Of those, $22$ completed, spanning $2.37 \times 10^{6}$ to
$8.52 \times 10^{7}$ states, and the remaining ones reached the larger cap as
well.

Counting each censored subtree at the cap gives a lower bound of
$7.9 \times 10^{11}$ states, $21\,838$ core-hours, or $65$ days on fourteen
cores. It is only a lower bound, and a soft one: at $\omega = 6$ the heaviest
subtree was $1.4 \times 10^{9}$ states, fourteen times this cap. If the censored
subtrees are five times the cap the figure becomes $206$ days; if they match the
$\omega = 6$ monster, $524$ days.

\subsection{Why the estimate never settled}
The revealing part is the sequence of estimates as the heavy subtrees completed,
in order of completion:

\begin{table}[h]
\centering
\begin{tabular}{lrr}
\toprule
heavy subtrees reported & estimate (days on 14 cores) & largest seen \\
\midrule
11 & 14 & $1.4 \times 10^{7}$ \\
19 & 21 & $4.0 \times 10^{7}$ \\
20 & 24 & $4.2 \times 10^{7}$ \\
22 & 36 & $8.5 \times 10^{7}$ \\
27 (5 censored) & $\ge 65$ & $> 10^{8}$ \\
\bottomrule
\end{tabular}
\end{table}

Throughout, the median of the heavy stratum barely moved, from
$6.2 \times 10^{6}$ to $8.5 \times 10^{6}$, while the mean more than quintupled;
by the end the five largest subtrees held $68\%$ of all sampled heavy states.
That is the signature of a distribution whose mean is not estimable from a few
dozen draws: each new maximum raises the average, so every reading understates
the last, and the sequence above has no visible limit. Our own reporting tracked
this upward, and it would have been easy --- and wrong --- to have stopped at any
point and announced the figure then in hand.

The methodological lesson is that stratifying removed one level of the problem
and revealed another: the heavy stratum is itself heavy-tailed. A sampling
estimate can bound such a search from below but cannot cost it, and the
progression of a bound that keeps rising is evidence about the distribution
rather than about the total.

\subsection{Cost is predictable before searching}\label{sec:predictor}
A search that cannot be finished whole can still be done in part, provided one
knows which parts are cheap. For this tree that turns out to be easy, and the
reason is Lemma~\ref{lem:range}: the admissible interval for the next prime runs
from $\lo/(2 - \lo + c/n_{\min})$ up to $r\hi/(2-\hi)$, and both ends are
products of fractions over the subtree's own components. Computing them costs
microseconds and requires no searching at all.

Against the $240$ sampled subtrees, the width of that interval predicts measured
cost with a log-log correlation of $\rho = 0.845$, states growing as
width$^{1.37}$; the cruder $r\hi/(2-\hi)$ alone gives $\rho = 0.824$. Binned by
prediction, the measured medians are monotone:

\begin{table}[h]
\centering
\begin{tabular}{lrrr}
\toprule
predicted interval width & subtrees & censored & median states \\
\midrule
$< 10^{3}$          & 23 & 0  & $1\,931$ \\
$10^{3}$--$10^{4}$  & 79 & 0  & $46\,891$ \\
$10^{4}$--$10^{5}$  & 48 & 13 & $1\,469\,617$ \\
$10^{5}$--$10^{6}$  & 42 & 27 & $> 2 \times 10^{6}$ \\
\bottomrule
\end{tabular}
\caption{Predicted against measured subtree cost. Every censored subtree falls in
the top two bins.}
\end{table}

Two uses follow. The first is practical: subtrees can be searched in ascending
predicted cost, so coverage accumulates quickly and the intractable states are
left identified rather than blocking the queue. Of the $80\,458$ depth-five
subtrees, $30\,259$ ($37.6\%$) are predicted below $10^{4}$; searching them in
that order took well under an hour, where an estimate from the \emph{average}
subtree cost would have suggested many hours. The second is diagnostic: the
predictor says where the remaining cost lives, which is how we identified the
configuration of Section~\ref{sec:fermat-alive}. A candidate new lemma can then be
tested against the expensive subtrees directly, in minutes, instead of by running
the search.

The predictor says nothing about one stratum. Where $\hi > 2$ ---
$14\,745$ subtrees, $18.3\%$ of the total --- Lemma~\ref{lem:range}(3) gives no
upper bound at all, so the interval has no width to measure. The figures above
exclude it, and the next subsection measures it instead.

\subsection{The unscored stratum is bimodal}\label{sec:hi2}
A state with $\hi > 2$ must contain an open component: were every component
fixed we would have $\hi = \lo > 2$, and the state would already be dead by
Lemma~\ref{lem:range}(4). That suggests such states are cheap, and the suggestion
is half right. Raising every open exponent together drives $\lo$ monotonically
towards $\hi > 2$, and in practice it crosses $2$ almost at once --- in a typical
case from $1.999994661$ to $2.000609256$ in a single step --- so that path dies
immediately. But the search also \emph{fixes} components, and a fixed component
contributes $h(p^{2a})$ in place of the larger limit $p/(p-1)$, so a fixed child
has a smaller $\hi$ which may fall back below $2$ and resume an ordinary search of
any cost. A cheap root does not imply a cheap subtree.

Measurement bears this out, and sharply. Of $400$ sampled subtrees from this
stratum, searched under a cap of $2 \times 10^{6}$ states:

\begin{table}[h]
\centering
\begin{tabular}{lr}
\toprule
quantile & states \\
\midrule
minimum          & $2$ \\
10th percentile  & $3$ \\
25th percentile  & $7$ \\
median           & $13$ \\
75th percentile  & $28$ \\
90th percentile  & $1\,760\,954$ \\
\bottomrule
\end{tabular}
\caption{The unscored stratum at $\omega = 7$. Three quarters of these subtrees
finish in under thirty states; a tenth exceed a million. The mean, $225\,584$, is
a property of the tail and describes almost none of them.}
\end{table}

So $318$ of $400$ ($80\%$) finish in under a thousand states while $39$
($10\%$) exceed two million. Extrapolated, the stratum costs at least
$3.3 \times 10^{9}$ states --- a lower bound, since the censored subtrees are
unbounded above --- and contributes roughly $1\,500$ subtrees to the expensive
part of the tree, an order of magnitude more than the $110$ with a finite
prediction above $10^{7}$. No solution and no unresolved state appeared.

The implementation therefore returns infinity for this stratum, which should be
read as \emph{unscored} rather than as a prediction of high cost: it keeps these
subtrees out of any band selected as cheap, which is the conservative choice. We
record that we nearly made the opposite error. Having verified the
$\lo$-crossing mechanism above, we changed the predictor to score them as
\emph{zero} --- the cheapest of all --- and caught it only because a probe
stalled with every worker busy. The mechanism was real but applied to one path
through each subtree rather than to the subtree, and a correct argument about a
part had been read as settling the whole.

\subsection{Partial verification}\label{sec:partial}
The sampling was not wasted as mathematics, even though it failed as
measurement: every subtree run to completion is genuine coverage of $\omega = 7$.
Using the predictor of Section~\ref{sec:predictor}, we then searched the cheap
part of the tree deliberately rather than incidentally. Of the $80\,458$
depth-five subtrees, the $30\,259$ with predicted width below $10^{4}$
($37.6\%$) were searched to completion in ascending predicted cost:
$4\,093\,731\,268$ states in $5.2$ hours on fourteen cores, none censored, with
no solution found and no state left unresolved. Together with the sampling, some
$5.2 \times 10^{9}$ states of $\omega = 7$ have been searched clean.

This must not be over-read. A partial search of a tree establishes nothing
whatever about the tree; $37.6\%$ of the subtrees is not $37.6\%$ of a proof, and
the states covered are a far smaller fraction of the total than the subtree count
suggests, since cost was the criterion for inclusion. What it provides is work
that need not be repeated, recorded subtree by subtree in \texttt{results/} with
each entry's predicted and actual cost.

That record also validates the predictor on a large sample rather than on the
$240$ it was fitted to. Its forecast for this band was $4.07 \times 10^{9}$
states against an actual $4.09 \times 10^{9}$, an error of $0.6\%$. Per subtree
the log-log correlation falls to $\rho = 0.688$, but that is restriction of
range rather than a failure: the band spans one decade of prediction where the
original fit spanned four. The aggregate accuracy is what matters for deciding
whether a band is affordable, and it is good.

A caution on a different reading of the same run, since we fell for it twice.
Because subtrees were searched cheapest-first, the progress estimate rose
steadily throughout --- an early extrapolation suggested well under an hour,
while the run took $5.2$. An elapsed-time projection is systematically
optimistic whenever work is deliberately ordered from cheap to dear, for exactly
the reason that the sampling estimates of Section~\ref{sec:omega7cost} were
optimistic.

\subsection{Conclusion of the costing}
A direct $\omega = 7$ search \emph{by the method of this paper} is not worth
starting. At best it is two months of continuous computation, at worst well over a
year, and the uncertainty cannot be reduced by more sampling of the same kind.
That is a verdict on this method, not on the problem: the missing ingredient is
structural, and Section~\ref{sec:theirlemmas} identifies it and finds the
computation affordable once it is supplied. The two conclusions are consistent and
should be read together --- enumeration of the two-free-prime leaves is hopeless,
and solving them is not.

% ---------------------------------------------------------------------------
\section{What the published solution does differently}\label{sec:theirlemmas}

Having reached the limits of our own method we read~\cite{tty} properly rather
than from its abstract. Three remarks follow: one correction to a comparison we
made, one identification of the exact gap in our approach, and a revised estimate
of what finishing $\omega = 7$ would cost.

\subsection{What we take on trust, and what we do not}\label{sec:trust}
A caution before any of it. \cite{tty} is an unrefereed preprint of August 2026,
and its central claim --- $\omega(n) \ge 8$ --- rests on a computation performed
once. Its own authors report that an earlier attempt at the same bound had, on
audit of its public code, a gap covering $35$ feasible
exponents~\cite{zemann}. So this is a setting with one documented silent failure
already, and the appropriate posture towards a second unreplicated computation is
not deference.

We therefore distinguish two things. Their \emph{lemmas} are mathematics and can
be checked here; we do so below rather than adopting them on authority, because a
prune implemented from an unverified lemma makes a search unsound in the one
direction that leaves no trace --- it discards solutions silently, which is
precisely the failure mode reported of~\cite{zemann}. Their \emph{computation} is
not something reading can check at all. Nothing in this section should be read as
confirming that $\omega(n) \ge 8$; what it establishes is that the tools are
sound and that repeating the computation independently is affordable.

We also record that our first description of these lemmas was taken from a summary
rather than from the paper, and stated the discriminant identity as though it had
been checked when it had not. It has since been derived here from scratch, as
below.

\subsection{The scan bounds agree}
We earlier reported our scan bound for the intermediate prime as
$8.59 \times 10^{9}$ against the $8 \times 10^{8}$ of~\cite{tty}, and read the
difference as an order of magnitude against us. That was a mistake of units.
Writing $\varepsilon = 2 - h(m)$ for the abundancy gap of the fixed part, that
paper gives the scan domain of a leaf as $(2/\varepsilon,\ 4/\varepsilon]$, of
width $2/\varepsilon$. Our own bounds give the same interval: by
Lemma~\ref{lem:range}(2) the lower end is
$\lo/(2 - \lo + c/n_{\min}) \approx 2/\varepsilon$, and by
Lemma~\ref{lem:range}(3) with $r = 2$ the upper end is
$r\hi/(2-\hi) \approx 4/\varepsilon$. Their figure is the width of that interval
at $\varepsilon \approx 2.5 \times 10^{-9}$; ours was an endpoint for a shallower
state. There is no discrepancy, and no order of magnitude to recover.

\subsection{The gap: two free primes}
The difference is not in the bounds but in what is done inside them. Our
Lemma~\ref{lem:lastprime} closes a leaf with \emph{one} unknown prime in closed
form, which is why $\omega = 6$ finished; with two unknown primes we have nothing
and must enumerate the first, which is the cliff of
Section~\ref{sec:fermat}. Their discriminant criterion closes exactly that case. We state it as we
derived it, for a leaf $n = m p^{2} q^{2}$ with $m$ fixed and both free primes at
exponent two. Writing $A = \sigma(m)$, $S = A(p^{2}+p+1)$ and $T = 2mp^{2}$, the
equation $\sigma(n) = 2n + c$ reads
$S(q^{2}+q+1) = Tq^{2} + c$, that is
\[
  (S - T)\,q^{2} + S\,q + (S - c) = 0 ,
\]
a quadratic in $q$. Hence
\[
  \Delta(p) = S^{2} - 4(S-T)(S-c)
\]
must be a perfect square, and $q$ is then determined as
$q = \bigl(S \pm \sqrt{\Delta(p)}\bigr) / \bigl(2(T-S)\bigr)$. At $c = 1$ the
discriminant is $S^{2} + 4(T-S)(S-1)$, which is the form~\cite{tty} gives, with
the same $S$ and $T$.

The derivation needs $S \ne T$, or the equation is linear in $q$ rather than
quadratic. That case cannot arise: $m$ is an odd square, so $A = \sigma(m)$ is a
product of terms $\sigma(p^{2a})$, each a sum of an odd number of odd numbers and
hence odd, and $\sigma(p^{2}) = p^{2}+p+1$ is odd, so $S$ is odd; while
$T = 2mp^{2}$ is even. This is the parity argument of Lemma~\ref{lem:parity}
again.

\begin{remark}
We checked this on genuine solutions rather than only deriving it. Taking every
$n = m^{2}p^{2}q^{2}$ in a small range for which $\sigma(n) = 2n + c$ holds with
$c > 0$ --- solutions exist for $c > 1$, which is what makes the test possible at
all --- all $838$ such leaves have $\Delta(p)$ a perfect square and recover the
correct $q$ exactly, with $q$ tested as an exact rational rather than by integer
division. The criterion is therefore sound independently of its source.
\end{remark}

A two-dimensional search becomes a scan over $p$ with a constant-time test, and a
candidate that fails leaves a certificate of its own rejection. Note the
restriction: the argument above is quadratic in $q$ only because $q$ has exponent
two, and $\sigma(q^{2b})$ has degree $2b$. Higher exponents need their third
lemma, which is presumably why it exists.

Their quadratic-residue sieve then filters $p$ cheaply beforehand by testing
$\Delta(p)$ against the quadratic residues of several small moduli. That one needs
no verification beyond its statement: a perfect square is a quadratic residue to
every modulus, so a single modulus witnessing otherwise rejects $p$ outright.
\subsection{The multilinear resolver, derived}\label{sec:resolver}
The third lemma cannot be skipped, because the discriminant argument needs the
last prime at exponent two. We therefore derived it as well. Substituting
$V = t^{2e}$ turns $\sigma(t^{2e})$ into $(tV-1)/(t-1)$, which is linear in $V$.

For a leaf $n = m t^{2e}$ with one unbounded exponent, $A = \sigma(m)$, the
equation $A(tV-1)/(t-1) = 2mV + c$ rearranges to
\[
  V = \frac{A + c(t-1)}{At - 2m(t-1)} ,
\]
so the exponent is determined outright by $t$: compute $V$ and test whether it is
a power of $t$. This is sharper than our Lemma~\ref{lem:oneexp}, which bounds the
exponent rather than fixing it. The denominator cannot vanish: $A$ and $t$ are odd
so $At$ is odd, while $2m(t-1)$ is even.

For two unbounded exponents, $n = m s^{2d} t^{2e}$ with $V_1 = s^{2d}$ and
$V_2 = t^{2e}$, put $B = A(sV_1 - 1)$ and $K = (s-1)(t-1)$. The equation is
bilinear in $(V_1, V_2)$ and gives
\[
  V_2 = \frac{B + Kc}{Bt - 2mK V_1} ,
\]
a M\"obius transformation of $V_1$, since $B$ is itself linear in $V_1$. As
$V_2 \ge t^{2}$ must hold, this confines $V_1$ to an explicit interval. Here parity
does not dispose of a vanishing denominator --- $B$ is even --- but it does not
need to: $B + Kc > 0$, so any solution forces $Bt - 2mKV_1 > 0$, and a
non-positive denominator is a rejection rather than an exception.

\begin{remark}
Both were tested on genuine solutions in the manner of
Section~\ref{sec:verification}: of $67$ leaves with one unbounded exponent and
$330$ with two, for which $\sigma(n) = 2n+c$ genuinely holds, every one recovers
its exponent exactly, with the arithmetic done over $\mathbb{Q}$ rather than by
integer division. The parity claims above were checked over $1\,329$ further
configurations, and the sign claim at every one of $1\,296$ genuine solutions.

An earlier version of these tests excluded the degenerate cases with a guard
rather than accounting for them, which would have concealed a failure there had
one existed. They are now shown impossible rather than merely unobserved, which
matters: a leaf closer that silently skips a case it cannot handle is the exact
mechanism by which an exhaustive search stops being exhaustive.
\end{remark}

All three lemmas are therefore available to us on our own authority rather than
on~\cite{tty}'s. Our Lemma~\ref{lem:oneexp} is superseded by the first half of
this one.

\subsection{Revised cost}
With those lemmas, \cite{tty} \emph{report} eliminating $79\,751\,212$ deep leaves
across $381$ stems in roughly $387$ core-hours. We cannot check that figure, and
record it as reported; but it is the right order for the arithmetic involved, and
set against our own lower bound of $21\,838$ core-hours for the same $\omega$ by
enumeration it is a factor of about fifty, which would bring the computation
within reach of the single desktop used here. Their deepest
leaves lie in a stem beginning $3, 5, 17, 137$, which is worth noting against the
$3, 5, 17, 257$ family that dominated our $\omega = 6$ computation
(Section~\ref{sec:fermat}); whether that reflects their exponent handling or a
genuinely different hardest configuration we have not determined.

\subsection{Why adopting their lemma would not have helped us}
\label{sec:wronglever}
We set out to implement the discriminant criterion inside the search of this
paper, expecting the factor of fifty to follow. It would not have. The criterion
makes a \emph{leaf} cheaper, and in our implementation leaves are almost free
already because almost nothing reaches one. Across the recorded
$\omega = 6$ and $\omega = 7$ work, $3.4 \times 10^{9}$ states produced $338$
exact leaf closures --- one in ten million --- and with the prunes of
Section~\ref{sec:lemmas} enabled, $\omega = 6$ performs \emph{no} leaf closures
at all in its $1.5 \times 10^{7}$ states. Essentially everything dies to the
abundancy bounds long before a leaf.

We implemented the criterion anyway, to measure rather than assume: with a
residue sieve over ten small moduli it tests candidate intermediate primes at
$1.46 \times 10^{6}$ per second against $4.4 \times 10^{5}$ for our existing
closure, a factor of $3.3$. The sieve itself is excellent --- it rejects
$99.96\%$ of non-squares, so only one candidate in $2\,800$ reaches an integer
square root, and it rejected none of $2 \times 10^{5}$ genuine squares. A full
scan of the $4 \times 10^{8}$ candidates in the widest window costs under five
minutes on one core. All of which optimises a thousandth of a percent of our
running time.

The conclusion is architectural, not mathematical. Our cost is in building and
pruning tree states with exact rational arithmetic; theirs is in scanning
integers. A factor of fifty is not available by adopting a lemma into this
structure --- it requires not building the tree over the last two primes at all,
and instead entering a flat integer scan once a stem and an exponent pattern are
fixed. That is a different program, and the lemma is the small part of it.

\subsection{What a verification would require}
The consequence is that $\omega = 7$ as an independent
verification is affordable in principle but needs a scanner rather than an
extension. That is the
right use of their work: borrowing a lemma we have proved for ourselves in order
to check a computation we have not. The bound has been computed once, and once
before that in a way its own successors report as having had a coverage
gap~\cite{zemann}; until a second computation exists, $\omega(n) \ge 8$ is a
claim resting on one unreplicated run.

% ---------------------------------------------------------------------------
\section{\texorpdfstring{$\omega = 7$}{omega = 7}: a complete search}\label{sec:omega7}
Section~\ref{sec:partial} ends with $\omega = 7$ out of reach: at least
$21\,838$ core-hours by the exact search, and probably several times that. This
section reports how it was brought within reach and then searched in full, with
no solution and nothing unresolved. With Section~\ref{sec:results}, that gives
$\omega(n) \ge 8$ for every quasiperfect number. This is the bound announced
by~\cite{tty}, found by an independent search with a different method and
different code.

\subsection{Where the cost was}
Profiling the dearest subtrees of the stratified sample (Section~\ref{sec:hi2})
showed one shape throughout: a node with two primes still to choose ($r = 2$)
and every exponent so far still open, such as
$3^{a} 5^{b} 17^{c} 257^{d} q_5^{2}$ followed by $q < q'$. That is the shape of
the node that closed $\omega = 17$ in Section~\ref{sec:omega17}, so the method
used there was generalised.

\subsection{The node search}\label{sec:qnode}
\texttt{qnode.py} handles any $r = 2$ node. Every component $x$ with exponent
$2b$ contributes $\log f(x, b) = L(x) - d_x(b)$, with $L(x) = -\log(1 - 1/x)$
and $d_x(b) = -\log(1 - x^{-(2b+1)})$, so $\sigma(N) = 2N + c$ becomes
\[
  L(q) + L(q') = K + \sum_x d_x(b_x) + \varepsilon, \qquad
  K = \log 2 - \sum_{\text{node}} L(p), \qquad
  0 < \varepsilon \le \frac{c}{2 N_{\min}}.
\]
The differences from the $\omega = 17$ program are these. The exponents of $q$
and $q'$ are components of the box like any other, since at $\omega = 7$ the last
two primes can be small. The term $\varepsilon$ is bounded explicitly, so the same
code runs for $c > 1$. That is what makes the test against real solutions
possible, below. The range of $q$ is walked in sieved segments rather than held
in memory. When nothing is left open but $q'$, its window is enumerated
directly. Node constants come from $80$-digit logarithms and all sums are in
double-double arithmetic, with the same slack of $10^{-35}$ on every comparison,
so the program can keep too much but never discards a solution. Every survivor
is settled by the original exact search. The splitting order, which always
splits the open component with the widest contribution, turned out to matter
as much as compilation: on the dearest subtrees it needs roughly $300$ times
fewer states than the round-robin order of the exact search.

\subsection{Validation}
Three checks, each independent of the run itself.
\begin{enumerate}
\item The $\omega \le 6$ reproduction: no solution, nothing unresolved, in
  $3.5$ seconds against $8$ minutes for the exact search
  (Section~\ref{sec:results}).
\item Real solutions. For every odd $k \le 3\,000\,001$, $n = k^{2}$ gives
  $c = \sigma(n) - 2n$, and wherever $c > 0$ the search for that $c$ and
  $\omega(n)$ must find $n$. All $42$ cases with $k \le 300\,001$ and
  $c \le 10^{6}$ were checked, and a systematic tenth ($941$ cases) of the $9\,401$
  with $c \le 10^{11}$. Every solution was found, and nothing was left unresolved.
\item The dearest sampled subtrees, which the exact search could not finish in two
  hours each, closed in under a second each, with no solution.
\end{enumerate}
The second check found two real bugs before the run. For very large $c$ the
allowance $\varepsilon$ is enormous, an exponential overflowed, and a lower bound
became NaN, so a range was silently dropped. And a wide final window was pushed
onto the work stack all at once and overflowed it. Neither can occur for
$c = 1$, but neither would have been caught by testing at $c = 1$ alone.

\subsection{The run}
The tree was split at depth $5$ into the same $80\,458$ subtrees as
Section~\ref{sec:partial}, and run on six cores from 17 October 2026. A uniform
sample of $244$ had suggested about $4$ core-hours in all. By the next morning
$80\,450$ subtrees had finished in about $38$ core-hours, and the last eight were
\[
  3^{10} 5^{8} 17^{8} 257^{8} q_5^{2} \cdots, \qquad
  q_5 \in \{65497, 65519, 65521, 65537, 65539, 65543, 65551, 65557\}.
\]
This is the Fermat obstruction of Section~\ref{sec:fermat} once more:
$65537 = F_4$, and with $3, 5, 17, 257$ and a fifth prime this close to it the
abundancy ceiling sits within about $10^{-9}$ of $2$. Neighbouring subtrees
showed the cost climbing steeply towards $F_4$ from above: $13$ minutes at
$q_5 = 65701$, $3.6$ hours at $65579$, $7.6$ hours at $65563$. Below $65480$ it
was a few seconds.

Within one of these subtrees the cost sits in a narrow band of $q$: for
$q_5 = 65701$, almost all of it lies between $2 \times 10^{7}$ and
$7 \times 10^{7}$. So the eight were re-run split by the range of $q$
(\texttt{qchunks.py}). Each was expanded to its $96$ $r = 2$ nodes in all, and
each node's range $[\,q_5 + 1, 2^{38})$ was cut at $600$ geometric points,
giving $57\,600$ pieces that tile the range exactly. Split this way, the test
subtree took the same time and the same number of states as in one piece.

\paragraph{Issue: a hand-off that did not scale.} Fourteen pieces came back
flagged unresolved, without solutions. Near $q \approx 4 \times 10^{9}$
everything but $q'$ is fixed, and $q'$ lay beyond the $2^{38}$ limit of the
compiled primality test. Each such case had been handed whole to the exact
search, millions of them, which overflowed the hand-off limit; the code then
correctly refused to report the pieces as closed. The window for $q'$ in this
situation is a few integers wide, so the compiled search now lists it itself up
to $4 \times 10^{18}$, widened by a relative margin for the precision of a
double, and leaves only the primality test to the exact code. One flagged piece
went from $10\,352$ seconds to $22$. The fourteen were set aside and redone.

\paragraph{Issue: the machine.} Twelve workers alongside another long job left
the computer barely responsive, and it had to be reset; the run resumed from its
records with nothing lost. It finished on four to eight workers.

\subsection{Result}
\begin{table}[h]
\centering
\begin{tabular}{lr}
\toprule
subtrees at depth 5 & $80\,458$ \\
of which split by $q$-range & $8$ ($57\,600$ pieces) \\
exact-search states & $1.17 \times 10^{9}$ \\
compiled-search states & $1.78 \times 10^{11}$ \\
processor time (records kept) & $61$ core-hours \\
solutions & none \\
unresolved states & none \\
subtrees or pieces missing & none \\
\bottomrule
\end{tabular}
\caption{The first complete $\omega = 7$ search (\texttt{finalize\_omega7.py},
\texttt{results/omega7\_final.json}). Coverage is checked against the task list
regenerated from the splitter, not against the records alone.}
\end{table}

\textbf{Status, 19 October 2026.} This run used a fixed comparison slack that
Section~\ref{sec:fpgap} shows to be insufficient for the small primes of
$\omega = 7$. It is being repeated in full with the corrected slack, and the
theorem below rests on the repeat. Its figures will replace these when it
finishes.

\begin{theorem}\label{thm:omega8}
Subject to the implementation, and to Cattaneo's theorem that a quasiperfect
number is an odd square, every quasiperfect number has at least eight distinct
prime factors.
\end{theorem}

This is the bound of~\cite{tty}, so it is not new. Its value is that it is
independent: a different decomposition, different code, and a method derived here
rather than taken from that paper. Proposition~\ref{prop:forced15} is subsumed by
it. The caveat is the same as everywhere in this paper and sharper here. Nearly
all the work was done by new code checked by tests, by planted and real
solutions, and by the exact search on every survivor, but by no one else.

% ---------------------------------------------------------------------------
\section{Soundness of the floating-point comparisons}\label{sec:fpsound}
The compiled searches (\texttt{qfast.py}, \texttt{qnode.py}) discard a region
only when a floating-point comparison says it cannot hold a solution. This
section bounds the rounding error of every such comparison, and shows that the
slack added to each one exceeds that bound. This argument was missing from the
first version of the paper, and writing it out exposed a real gap, described in
Section~\ref{sec:fpgap}.

\subsection{Model}
Arithmetic is IEEE binary64 with round-to-nearest, $u = 2^{-53}$, with no fused
multiply-add. A double-double value is a pair $(h, \ell)$ with $|\ell| \le u|h|$,
representing $h + \ell$. Every quantity used lies far from overflow. Underflow
can occur only in powers $x^{-(2b+1)}$ below $10^{-290}$, and those enter only as
upper bounds on $d_x(b)$, where a value below $10^{-290}$ is covered by the
absolute slack. We use three standard facts. \textsc{TwoSum}$(a, b)$ returns
$s, e$ with $s + e = a + b$ exactly (Knuth). \textsc{FastTwoSum}$(a, b)$ does the
same when $|a| \ge |b|$. Dekker's \textsc{TwoProd} returns $p + e = ab$ exactly
when nothing underflows.

\subsection{The primitives}
\begin{lemma}\label{lem:ddadd}
For double-double inputs $a, b$, the addition in \texttt{dd\_add} returns $c$
with $|c - (a + b)| \le 20u^{2}(|a| + |b|)$.
\end{lemma}
\begin{proof}
\textsc{TwoSum} gives $s + e_0 = a_h + b_h$ exactly, with $|e_0| \le u|s|$. The
two further additions of $a_\ell$ and $b_\ell$ to $e_0$ incur at most
$2u(1+u)(|e_0| + |a_\ell| + |b_\ell|) \le 2u^{2}(1+u)(2+u)(|a_h| + |b_h|)$, since
$|s| \le (1+u)(|a_h| + |b_h|)$. If the final \textsc{FastTwoSum} has its
precondition it is exact. If it does not, $|s|$ is smaller than the corrected
term $\hat e$, and its error is at most $2u(|s| + |\hat e|) \le 4u|\hat e|$,
which is again $O(u^{2})(|a| + |b|)$ by the previous bound. Adding the two cases
gives the constant $20$, which is not tight.
\end{proof}

\begin{lemma}\label{lem:ddseries}
The routine computing $-\log(1 - x)$ for double-double $0 < x \le 1/7$, by summing
$x^{k}/k$ until a term falls below $10^{-48}$, has error at most
$1200u^{2}\,|{-\log(1-x)}| + 2 \times 10^{-48}$. The same bound covers
$L(x) = -\log(1 - 1/x)$ for $x \ge 7$ and $d_x(b)$ for all $b \ge 1$.
\end{lemma}
\begin{proof}
Each power $x^{k}$ is formed by $k - 1$ double-double products, each with relative
error at most $8u^{2}$, and divided by $k$ with relative error at most $2u^{2}$,
so term $k$ carries relative error at most $8ku^{2}$. Weighted by the terms, that
is at most $8u^{2} \sum_{k} x^{k} \le 10u^{2}\,(-\log(1-x))$ for $x \le 1/7$.
Summing at most $60$ terms costs, by Lemma~\ref{lem:ddadd}, at most
$60 \times 20u^{2}$ times the partial sums, all bounded by the total. Truncation
drops a tail below $10^{-48}/(1-x) \le 1.2 \times 10^{-48}$. For $d_x(b)$ the
argument is $x^{-(2b+1)}$, formed by $2b$ products, which adds at most $16bu^{2}$
relative error to a term that is tiny whenever $b$ is large. The bound stands
with margin at the exponents that occur.
\end{proof}

The node constants, $K$ and every $d_p(a)$ of a node prime, are computed
separately as $80$-digit decimals from exact rationals and then split into
$(h, \ell)$. Each is therefore within $2u^{2}$ of its true value, relatively.

\subsection{The comparisons}
Every pruning comparison has the form $\sum_i \pm T_i \ \lessgtr\ \pm\sigma$, where
the $T_i$ are $K$, the $d$-terms of the box, and $L(q)$ and $L(q')$ where present.
The sum takes at most $16$ additions for $\omega \le 7$ and at most $20$ for the
$\omega = 17$ node. Combining Lemmas~\ref{lem:ddadd} and~\ref{lem:ddseries}, the
computed sum is within
\[
  E \;=\; (16 \times 20 + 1200 + 2)\,u^{2} \sum_i |T_i| \;+\; 4 \times 10^{-48}
  \;<\; 2048\,u^{2} \sum_i |T_i| \;+\; 10^{-47}
\]
of the exact sum. The double-double result is then rounded to one double, a
relative change of at most $u$. That preserves the sign, and costs at most $u$
times the slack, which the margin in $2048$ absorbs. Integer bounds for the next
prime are derived from these values through $\exp$, $\log$, $\sqrt{\ }$ and
reciprocals. That adds a relative error of at most $10u$ in $q$, which is covered
by the margins of $\pm 2$ and the relative $\pm 4 \times 10^{-15}$ used there.
The allowance $\varepsilon \le c/(2N_{\min})$ is computed from a sum of logarithms,
with relative error below $10^{-11}$, and multiplied by $1 + 10^{-6}$.

\begin{proposition}\label{prop:fpsound}
If every comparison uses slack at least $10^{-35} + 2048u^{2} \sum_i |T_i|$, no
comparison of the compiled search discards a region containing a solution.
\end{proposition}
\begin{proof}
The slack exceeds $E$ in every comparison, so a comparison that discards a region
would also discard it in exact arithmetic. Each such exact test is a consequence
of the equation $L(q) + L(q') = K + \sum d + \varepsilon$ restricted to the box.
\end{proof}

\paragraph{The $\omega = 17$ node.} \texttt{qfast.py} uses a fixed slack of
$10^{-35}$. Every term there has size at most about $2 \times 10^{-9}$, so
$\sum |T_i| \le 10^{-8}$ and $E < 2.6 \times 10^{-37}$, a fortieth of the slack.
Its series are short and explicit: $\log f(x)$ to four terms in
$x^{-1} \le 1.8 \times 10^{-9}$, and $d_p(a)$ and $K$ from rationals. Their
truncation is below $10^{-44}$. So Proposition~\ref{prop:fpsound} applies as the
code stands, and the $\omega \ge 18$ result does not depend on the gap below.

\subsection{The gap, and the re-run}\label{sec:fpgap}
In \texttt{qnode.py} the terms are not small. For a node such as
$3^{a} 5^{b} 17^{c} 257^{d} q_5^{2}$ followed by $q$ and $q'$, the prime $q$ can be
as small as $q_5 + 1$, and at the smallest nodes $L(q)$ approaches $0.15$. There
$2048u^{2}\sum|T_i|$ is about $10^{-30}$, and the measured error of $L(7)$ alone
is $2.9 \times 10^{-33}$, against a fixed slack of $10^{-35}$. A true solution
sits exactly on the boundary of its box, with margin $\varepsilon$, about
$1/(2N)$, far below either figure. So the first $\omega = 7$ run (17--19 October
2026) did not satisfy Proposition~\ref{prop:fpsound}. A solution could in
principle have been discarded by rounding, though nothing suggests one was.

\texttt{qnode.py} now adds $2048u^{2}\sum|T_i|$ to the slack of every comparison,
and $\omega = 7$ has been run again in full on the corrected code; the result
reported in Section~\ref{sec:omega7} is the re-run. The records of the first run
are kept, marked \texttt{v1}.

\paragraph{Measurement.} \texttt{test\_dd\_error.py} compares the kernel's $L$,
$d$ and chains of the comparison shape against $90$-digit references, on $4\,000$
random cases with $x$ from $7$ to $10^{12}$. The worst errors beyond the
truncation term are $3.9u^{2}$, $9.2u^{2}$ and $4.6u^{2}$ times the magnitude. That
is more than two hundred times inside the bound, which is deliberately loose.

\section{How the search was developed}\label{sec:history}

The method of Section~\ref{sec:lemmas} was not designed in one pass. This
section records, in order, the problems met while building and running it: what
went wrong, why, and what was changed. Several of the final design choices only
make sense in the light of these problems, and two of them were errors of ours
that the testing caught.

\subsection{Choice of approach}
Two routes were considered for the last prime of a branch. One enumerates
candidate primes up to the bound of Lemma~\ref{lem:range}(3). The other solves
for the prime directly. Enumeration fails exactly where the problem is hardest:
when the abundancy of the fixed part is very close to $2$, the bound on the last
prime grows like $1/(2 - h(m))$ and becomes astronomically large. We therefore
chose from the start to close such leaves algebraically
(Lemma~\ref{lem:lastprime}), and to handle unbounded exponents by carrying them
as intervals rather than enumerating them. The price of intervals is that a leaf
can have an unknown exponent as well as, or instead of, an unknown prime; that
is what Lemma~\ref{lem:oneexp} and the parity argument of
Lemma~\ref{lem:parity} are for.

The decision to solve the general equation $\sigma(n) = 2n + c$, rather than
only $c = 1$, was made for testing. For $c = 1$ there are no known solutions, so
a search that wrongly prunes everything would look identical to a correct one.
For other values of $c$ solutions exist and can be found by brute force, so a
wrong prune shows up as a missed solution.

\subsection{Issue 1: an incomplete root finder}
The first version of the last-prime solver bisected $P_a$ only on the stretch
beyond its maximum, and proved that no root lay on the rising stretch before it
by a size inequality. That inequality holds comfortably when $c$ is small, but
not in general. The first round of planted tests, whose values of $c$ are huge,
produced many cases where the inequality failed. The search flagged each of
these rather than skipping them, so no case was silently lost. The fix was the
exact method now described after Lemma~\ref{lem:lastprime}: locate the turning
points of $P_a$ through the single sign change of $P_a''$, and bisect every
monotone stretch.

The two versions have been re-run against each other on the same $200$ planted
last-prime cases, from the same random seed
(Section~\ref{sec:variants} explains how the faulty version was recovered). The
original recovered $45$ of the $200$ planted solutions and flagged $190$ cases
as unresolved; the corrected version recovers all $200$ and flags none. So the
defect was severe, and it was visible rather than silent: a run carrying it could
never have been mistaken for a proof.

\subsection{Issue 2: a case that cannot occur for $c = 1$}
Lemma~\ref{lem:range}(4) shows that for $c = 1$ a branch whose fixed part is
already abundant is dead. For larger $c$ it is not, and the first brute-force
comparison flagged one such state (with $c = 189\,651$). The lemma in fact gives
a finite bound, $q < c/(n_{\min}\lo)$, on the next prime, and the search now
enumerates that range. For $c = 1$ the range is empty, so this changes nothing
for the real problem; it matters only for the tests.

\subsection{Issue 3: a test that could not finish}
The first planted-solution test ran the whole search for each planted $c$. It
stalled: planted values of $c$ are comparable to $n$, which makes the range of
Issue~2 enormous. This was a flaw in the test, not in the search, since that
range is empty for $c = 1$. The test was redesigned to exercise the two closing
lemmas directly on planted data, which is where brute force cannot reach. The
search as a whole remains covered by the brute-force comparison.

\subsection{Issue 4: five primes too slow}
The first run with $\omega = 5$ was stopped after more than half an hour. By
sampling the states being visited, we found that the search, with one prime left
to choose and several exponents still open, was enumerating the last prime over
ranges of tens of thousands, each candidate carrying its own refinement subtree.
The admissible range for the last prime lies between
$\lo/(2-\lo)$ and $2/(2-\hi)$, and its width tends to $1$ as the open intervals
tighten. Refining the open exponents first therefore collapses the range to one
or two candidates. With that change $\omega = 5$ took three seconds.

\subsection{Issue 5: the same idea, applied too widely}
Applying the same rule with two or more primes left made $\omega = 5$ ten times
slower (33\,s against 3\,s). With several primes still to come, the admissible
range for the next prime does not shrink to a point as the exponents tighten; it
only settles to a ratio of about $r$ between its ends. So early refinement
spends effort without buying much pruning. The change was reverted, and the
rule kept for $r = 1$ only (but see Issue~8).

\subsection{Issue 6: a long run lost}
The first $\omega = 6$ run was a single process, and it died when the working
session that launched it ended, after about four hours and $2.6 \times 10^8$
nodes. It had not finished, and it gave no indication of how far it had got. It
was replaced by a driver that cuts the tree at a fixed depth into independent
subtrees, runs them in parallel, and records each finished subtree, so that an
interrupted run resumes where it stopped. Because each subtree is an ordinary
call of the same search, the cut does not change what is searched; this was
confirmed by reproducing the serial node count for $\omega = 5$ exactly. The
run was also started as a process independent of the session.

\subsection{Issue 7: one subtree holding most of the work}
Splitting at four primes gave 358 subtrees, most finishing in seconds and a few
taking up to two hours. One, beginning $3, 5, 17, 257$, outlasted all the rest
by hours. The reason is arithmetical: these are the Fermat primes
$2^{2^i} + 1$, and
\[
  \frac{3}{2}\cdot\frac{5}{4}\cdot\frac{17}{16}\cdot\frac{257}{256}
  = \frac{65535}{32768} = 2 - 2^{-15}.
\]
So four primes can bring the abundancy within $2^{-15}$ of $2$, leaving the
remaining two primes free to be very large, and the search must rule out a very
long range of them. This is the kind of ``deep leaf'' that Toyohara, Tao and
Yao~\cite{tty} identify at $\omega = 7$; it appears here already at $\omega = 6$.
Within this subtree the heaviest part is near the fifth prime $65537$, the next
Fermat prime, where the product comes within $2^{-31}$ of $2$.

To use the cores left idle, this subtree was itself split by its fifth prime,
giving $12\,195$ further pieces run in parallel. A splitting tool was written
for this and checked against an already finished subtree, whose node count it
reproduced exactly ($200\,601$).

\subsection{Issue 8: repeated work, and a loop}
The split run showed why the Fermat subtree was slow: for each of its twelve
thousand fifth primes, the search refined the exponents of $3$, $5$, $17$ and
$257$ afresh, repeating the same work thousands of times. The remedy is to refine
those exponents once, before choosing the fifth prime, but only when the next
prime's range is huge, so that the ordinary cases of Issue~5 are unaffected.
Most exponent choices then pull the product well away from $2$, and those
branches die at once.

The first version of this rule tested only $\hi$, and it does not terminate. For
an open exponent $\hi$ uses the limit $p/(p-1)$, which does not change as the
exponent's lower bound rises. Refining an open component yields one child with it
fixed and one with it still open, and on the all-open path $\hi$, $r$ and the set
of open components are all unchanged, so the rule fires again, forever. Nothing
else stops that path: $\hi < 2$ rules out the $\hi \ge 2$ refinement, and
$\lo \nearrow \hi < 2$ means Lemma~\ref{lem:range}(4) never applies either. The
corrected rule also uses $\lo$, which does rise with the exponents: it refines
only while the upper end of the next prime's range exceeds $4r$ times the lower
end. Since $\lo \to \hi$, the lower end rises towards $\hi/(2-\hi)$ while the
upper end stays at $r\hi/(2-\hi)$, so the condition eventually fails and
refinement stops. The rule changes the order of the search, never what it covers.

We first suspected this only because a timing test failed to produce output, which
is weak evidence: that subtree is genuinely slow. The claim has since been tested
directly on the four-Fermat-prime state, using the exponent floor of an open
component as a tripwire, since Lemma~\ref{lem:parity} bounds that floor in any
terminating search while a runaway path drives it up without limit. The faulty
version passes $2 \cdot 400$ --- a factor $3^{800}$ --- after only $256\,506$
nodes. The corrected version searches $10^{6}$ nodes of the same state with its
deepest open exponent never exceeding $2 \cdot 15$. Termination itself is not
observable here, because the subtree takes hours either way; that is what
Proposition~\ref{prop:terminate} is for. Constructing this test corrected a
mistake in our first attempt at one. Running the faulty version on $\omega = 5$ shows no
difference at all, because the rule never fires there: with two primes still to
choose, only three are fixed, and the largest $\hi$ available is
$\tfrac32\cdot\tfrac54\cdot\tfrac{17}{16} = \tfrac{255}{128}$, which is not
within $2r/5000$ of $2$. The rule needs four Fermat primes to fire, so the defect
is invisible below $\omega = 6$ --- the very case it was introduced to
accelerate. A speed heuristic added for a hard case can therefore be untestable
on every easy case, which is an argument for proving termination
(Proposition~\ref{prop:terminate}) rather than inferring it from runs.

\subsection{Issue 9: which run benefits}
Three attempts on the Fermat subtree were running concurrently when the corrected
rule was written: the original single process, the twelve-thousand-piece split of
Issue~7, and a fresh whole-subtree run on the new code. It is worth being precise
about which of them the correction helps. Every piece of the split has its fifth
prime fixed and one prime left to choose, so $r = 1$ throughout, and the new rule
applies only when $r \ge 2$. The split therefore gains nothing from it, and
restarting it on the newer code would be wasted effort. Only the whole-subtree
run, which passes through the $r = 2$ states, is affected.

\subsection{Issue 10: looking for a result instead of more compute}
With $\omega = 6$ settled and $\omega = 7$ apparently years away, the natural
move was to price the larger computation. The better move, which came only when
we stopped costing hardware and asked what was known about the equation itself,
was Lemma~\ref{lem:qr}. It is elementary, it took an afternoon to find, prove and
test, and it is worth about a factor of $100$ at $\omega = 6$ --- far more than
any scheduling or parallelism change in this project, and more than the machine
could have delivered by running longer.

Two details of the implementation were not obvious. The first attempt factored
$\sigma(p^{2a})$ hard, with Pollard rho on numbers up to $90$ digits, and made
$\omega = 5$ \emph{slower} than no prune at all: the factoring cost exceeded the
search it saved. Capping the effort at $25$ digits gave the $12.7$-fold gain.
The prune is worth having only where it is cheap, which is exactly where the
tree is dense. The second is that the cheap Jacobi test on the whole of
$\sigma(p^{2a})$ is not equivalent to the real condition, because two violating
primes cancel; $\sigma(257^{2}) = 61 \cdot 1087$ is such a case, and treating
the Jacobi test as decisive would have made the prune unsound in the permissive
direction.

\subsection{Issue 11: the literature, searched too late}
Lemma~\ref{lem:qr} was derived here from scratch, proved, tested, implemented and
measured --- and only then found in the literature, under a name, already
formalised in Lean~4 by~\cite{ualbf}. The derivation stands as an independent one
and is evidence of a kind, as argued in Section~\ref{sec:lemmas}; the failure was
not in the mathematics but in the order of work. We had suspected the result was
known and said so, but a suspicion recorded as a hedge is not the same as having
looked. A single search on the terms of the condition would have found it in
minutes.

The more costly omission was the second one. The same search would have turned up
the conditional bounds of~\cite{prasad}, $\omega(n) \ge 15$ when
$\gcd(n,15) = 1$, which bear directly on which parts of our own search are worth
running at all, since they cover precisely the configurations we go on to measure
as cheap (Section~\ref{sec:fermat}). Knowing that at the outset would have
redirected effort rather than merely saved a derivation.

Both were found by accident, while checking the page numbers of a 1951 citation.
One corollary is that verifying citations properly is not clerical work: it was
the most productive half hour of the project.

The other corollary needs stating carefully, because the obvious lesson --- read
first --- is not quite the right one. This work began from very little: the
statement of the problem, the two Hagis--Cohen bounds, Cattaneo's theorem, and
the abstracts of two 2026 preprints. Everything else was derived here. For the
central claim that is a virtue rather than an accident to apologise for. The
paper's value lies in reproducing Theorem~\ref{thm:main} \emph{independently},
and we never read the 1982 argument; had we done so before writing the search we
might have reimplemented its tree without noticing, and the reproduction would
corroborate much less. The resulting program is visibly unlike what that paper
describes --- unresolved exponents carried as intervals, leaves closed
algebraically rather than enumerated --- and that difference is exactly what makes
the agreement worth something. Working out the machinery before reading about it
also produced a tool we understood well enough to extend, which is how the cost
predictor of Section~\ref{sec:predictor} came about.

What that argument does \emph{not} excuse is ignorance of the surrounding
literature. The Legendre--Cattaneo filter and the conditional bounds
of~\cite{prasad} are different results from the one being reproduced; reading them
could not have compromised any independence, and not reading them cost real
effort. The distinction is the lesson: derive the method yourself if independence
is the point, and read everything that is not the method immediately. We did the
first by luck and the second too late.

\subsection{Lessons}
Five things recur. First, every issue above was found because the program reports
states it cannot close instead of discarding them, and because it was tested on
problems with known answers; Issue~1 flagged $190$ of $200$ planted cases rather
than quietly returning nothing. Second, the costly failures were about the order
of work rather than its correctness, and the useful diagnostic each time was to
sample which states the search was visiting. Third, a rule about when to refine
must be checked for termination separately from its effect on speed: Issue~8's
first version satisfied Proposition~\ref{prop:complete} exactly and would never
have finished. Fourth, and least expected: an optimisation aimed at the hardest
case may have no effect on any case small enough to test quickly, so its defects
cannot be found by regression testing at all. That is an argument for proving the
properties one needs rather than observing them. Fifth, and the most useful
lesson of the project: we spent far longer on parallelism, scheduling and
subtree decomposition than on the arithmetic, and the arithmetic
(Lemma~\ref{lem:qr}) was worth more than all of it combined. Engineering effort
on a search of this kind is easy to spend and hard to notice spending; a single
congruence condition beat every such change by two orders of magnitude.

% ---------------------------------------------------------------------------
% ---------------------------------------------------------------------------
\section{Reproducibility and availability}\label{sec:repro}

\subsection{Environment}
All computations were run on one desktop machine: an Intel Core i5-13400F
(13th generation, 16 logical cores) under Windows~11, with CPython 3.13.14.
The exact search (\texttt{qp.py}) uses only the standard library:
\texttt{fractions.Fraction} for rational bounds and Python's arbitrary-precision
integers throughout, with no floating-point arithmetic. The compiled searches of
Sections~\ref{sec:omega17} and~\ref{sec:omega7} (\texttt{qfast.py},
\texttt{qnode.py}) are different. They use Numba~0.68 and NumPy, and double-double
floating point with an explicit slack on every comparison. They only discard
regions, and every region they keep is settled by the exact search, so a float
error could cost time but could not create a false ``no solution'' unless it
exceeded the slack; the margin is about six orders of magnitude. No
computer-algebra system is used. Typesetting used Tectonic 0.17.0.

\subsection{Running it}
\begin{itemize}
\item \texttt{python qp.py K C} eliminates $\omega = 1, \dots, K$ for
  $\sigma(n) = 2n + C$ in a single process.
\item \texttt{python run\_parallel.py K DEPTH WORKERS} cuts the $\omega = K$
  tree at \texttt{DEPTH} primes into independent subtrees and runs them in
  parallel, appending one record per finished subtree to
  \texttt{results/omegaK\_dDEPTH.jsonl}. Re-running the same command resumes
  from that file.
\item \texttt{python split\_task.py K DEPTH TASK WORKERS} splits one subtree
  further, for the case of Issue~7, and records a single merged record for the
  original subtree so the main run treats it as complete.
\item \texttt{python test\_qp.py} and \texttt{python test\_leaves.py} are the
  two verification suites of Section~\ref{sec:verification}.
\item \texttt{python run\_qnode.py 7 5 WORKERS} runs the complete $\omega = 7$
  search of Section~\ref{sec:omega7}, with every $r = 2$ node done by the compiled
  kernel; \texttt{python qchunks.py OUT CHUNKS WORKERS @TASKS} splits dear
  subtrees by $q$-range; \texttt{python finalize\_omega7.py} checks coverage and
  writes the verdict.
\item \texttt{python test\_qnode.py w 6} reproduces $\omega \le 6$ with the
  kernel, and \texttt{python test\_qnode.py b K CMAX [STEP]} checks it against every
  brute-force solution of $\sigma(n) = 2n + c$ with $n = k^{2}$, $k \le K$.
\item \texttt{python run\_qfast.py results/cond\_3-5\_17.jsonl} closes the last
  $\omega = 17$ node; \texttt{test\_qfast.py} is its planted-solution test and
  \texttt{spot\_qrange.py} its spot check by the exact code.
\item \texttt{python audit\_results.py K DEPTH} audits a finished run for
  complete coverage, as described in Section~\ref{sec:results}, without
  repeating the search.
\item \texttt{python paper/reconstruct\_variants.py SESSION.jsonl} rebuilds the
  source variants of Section~\ref{sec:variants} and checks them, and
  \texttt{python paper/verify\_variants.py} runs each faulty variant against the
  defect the paper attributes to it, and its replacement against the same test.
\end{itemize}
A run constitutes a proof only if it reports \texttt{unresolved = 0}. States the
search cannot close are recorded, never skipped, so this is a meaningful check
rather than a formality.

\subsection{Data}
Each record in a results file carries the subtree it covers, its node count,
the number of exact leaf checks performed, the solutions found, the unresolved
states, and its wall-clock time. The union of the recorded subtrees is exactly
the tree of Proposition~\ref{prop:complete}, which can be re-derived from the
code without re-running the search, so the claim of complete coverage is
checkable independently of trusting our timings.

\subsection{Source variants}\label{sec:variants}
Section~\ref{sec:history} describes versions of the program that were written,
tested, found wrong and corrected. Two of those were corrected before being
committed to version control, so they survive in no commit, and the development
environment's own file backups retain nothing for files created during the
session. They are nonetheless part of the record: a reader checking Issue~1 or
Issue~8 should be able to run the faulty version and see it fail.

The session transcript records the complete text of every file written and the
exact before-and-after strings of every edit. Replaying those in order
reconstructs each intermediate version exactly, and the replay is self-checking:
where the final reconstructed state of a file matches the file now on disk byte
for byte, every earlier state in that chain is faithful too, since each step is
an exact textual substitution. The accompanying
\texttt{paper/variants/} directory holds the result, one file per version, with
a manifest giving each version's timestamp, SHA-1 and role; files whose chains
do not verify are marked as such and should not be relied on. Of the eight
chains recovered, seven verify; the one that does not is a test script that was
once edited through the shell, so the transcript holds the command but not the
resulting text.

The companion appendix \emph{Source code} (\texttt{code.pdf}) lists every
current source file in full, followed by the change that introduced and the
change that removed each defect of Section~\ref{sec:history}, so the claims made
there can be read against the code rather than taken on trust.

\subsection{What was not preserved}
The first $\omega = 6$ attempt (Issue~6) was a single process that was killed
after roughly four hours and $2.6 \times 10^{8}$ nodes without completing. Its
per-subtree detail was never written to disk and is lost; only the node count
and the absence of any flagged state are known, from its progress log. Nothing
of mathematical consequence depends on it, since the region it covered was
searched again by the parallel run. The wall-clock times of the three concurrent
attempts on the Fermat subtree of Issue~7 are likewise not a clean benchmark,
because they shared a machine.

% ---------------------------------------------------------------------------
\section{Further work}\label{sec:further}

\begin{itemize}
\item A direct $\omega = 7$ search is not advisable on one machine, and the
  attempt to cost it is instructive in itself
  (Section~\ref{sec:omega7cost}).
\item Attempt $\omega = 7$ as an independent check of~\cite{tty}. Toyohara, Tao
  and Yao report ``deep leaves'' at $\omega = 7$ on which pure enumeration does
  not terminate; whether Lemmas~\ref{lem:oneexp} and~\ref{lem:lastprime} are
  enough to close them is the main open question for this approach. The
  $\omega = 6$ run gives a concrete reason for caution and a concrete reason for
  hope. Caution, because the analogous obstruction appeared there too, in the
  four-Fermat-prime subtree, and consumed $70\%$ of the computation by itself.
  Hope, because it did close, by the same two lemmas, with no state left
  unresolved. At $\omega = 7$ the corresponding states will include a fifth
  Fermat prime, $65537$, bringing $\prod p/(p-1)$ within $2^{-31}$ of $2$; the
  admissible window for the remaining primes widens accordingly, and the cost
  is likely to be several orders of magnitude greater rather than merely larger.
\item A lower bound on $n$ itself: the half of Hagis and Cohen not reproduced
  here, standing at $10^{45}$ since~\cite{alekseyev}. We previously suggested in
  this section that the factor-chain method of Ochem and Rao~\cite{ochemrao}
  would transfer, and that the gap against the $10^{1500}$ known for odd perfect
  numbers indicated neglect. That was wrong, and the reason is instructive. For
  an odd perfect number $\sigma(n) = 2n$, so $\sigma(p^{k}) \mid 2n$ and the odd
  part of $\sigma(p^{k})$ \emph{divides} $n$; each prime power thus forces
  further prime factors of $n$, and the chain propagates. For a quasiperfect
  number $\sigma(n) = 2n+1$ is coprime to $n$, so by Lemma~\ref{lem:coprime} the
  primes of $\sigma(p^{k})$ are precisely those that \emph{cannot} divide $n$.
  The chain does not propagate at all. The disparity between $10^{45}$ and
  $10^{1500}$ therefore reflects a structural obstacle rather than a lack of
  attention, and a size bound here needs a different instrument --- such as the
  general solver for $a\sigma(n) = bn + c$ of~\cite{alekseyev} --- rather than a
  transplanted technique.
\item Removing the dependence on Cattaneo's theorem, by running the same search
  over $n = 2^{a}m^{2}$ with $a \ge 1$.
\end{itemize}

% ---------------------------------------------------------------------------
\section{Conclusion}\label{sec:conclusion}

It is worth stating plainly what this work does and does not establish, since the
honest total is smaller than the length of the paper might imply.

\textbf{Established.} No quasiperfect number has six or fewer distinct prime
factors (Theorem~\ref{thm:main}). The search is exhaustive by
Proposition~\ref{prop:complete}, terminates by
Proposition~\ref{prop:terminate}, leaves no state unresolved, and its coverage was
audited separately from the run. The result was obtained three times by routes
that differ in process structure, decomposition and pruning, twice agreeing to the
exact state count. It is conditional on Cattaneo's theorem~\cite{cattaneo}, which
we cite rather than reprove.

\textbf{Independently derived, not new.} Lemma~\ref{lem:qr} was proved, tested and
measured here before being found in the literature as the Legendre--Cattaneo
filter~\cite{ualbf}; the independence is documented and checkable, and the
convergence corroborates the implementation, but the result is not ours.

\textbf{Independent in the sense that matters.} The method was built from the
problem statement rather than from~\cite{hagiscohen}, which we have not read.
That is deliberate in effect if not in planning: a reproduction whose author has
studied the original argument corroborates far less, and the program here is
visibly unlike the one that paper describes. The same innocence about the
\emph{rest} of the literature, by contrast, cost effort and gained nothing
(Section~\ref{sec:history}, Issue 11).

\textbf{Not new.} Theorem~\ref{thm:main} is one of the two bounds of Hagis and
Cohen~\cite{hagiscohen}, from 1982. Their second bound, $n > 10^{35}$, since
raised to $10^{45}$~\cite{alekseyev}, is untouched here, and the factor-chain
technique that carried odd perfect numbers to $10^{1500}$ does not transfer
(Section~\ref{sec:further}). We also claim no originality for
Lemma~\ref{lem:qr}, which is elementary and may well be the
``quadratic-residue sieve'' of~\cite{tty}; nor for Lemma~\ref{lem:fermat}, which
follows from a classical identity. The value of the work, such as it is, lies in
the independence of the implementation from the original, not in the statements.

\textbf{Not achieved.} $\omega = 7$ is not settled and we do not recommend
attempting it this way (Section~\ref{sec:omega7cost}). Roughly $10^{9}$ states of
it have been searched with nothing found, but a partial search of a tree
establishes nothing whatever about the tree, and that figure should not be read as
evidence either way.

\textbf{Perhaps of wider use.} Three things here might be worth more than the
theorem. Lemma~\ref{lem:qr} and its corollary are worth a factor of $130$ at
$\omega = 6$ and were
found by asking what was true of the equation rather than how to compute faster
--- after we had spent considerably longer on parallelism than on arithmetic.
The cost of a subtree is predictable before searching it
(Section~\ref{sec:predictor}), which makes the tractable part of an intractable
search accessible and localises the rest. And the attempt to cost $\omega = 7$
failed in an instructive way (Section~\ref{sec:omega7cost}): a sequence of
sampling estimates that rises monotonically as a heavy tail arrives is evidence
about the distribution rather than about the total, and we recorded our own
drift up that sequence because it would have been easy to stop anywhere along it
and publish the figure then in hand.

\textbf{What would actually advance matters.} Not hardware. The cost at
$\omega = 7$ is concentrated in states where $3, 5, 17, 257$ occur at exponents
holding $\prod p/(p-1)$ within about $10^{-5}$ of $2$, with two primes still to
choose, and the lever is the scan bound for the next prime: ours stands at
$8.59 \times 10^{9}$, while~\cite{tty} reports $8 \times 10^{8}$. An order of
magnitude there is worth more than any machine available to us. Failing that, the
cheaper and still useful exercise is to implement the lemmas of~\cite{tty} and
check their computation, which has been performed once and, on their own account,
once before that with a gap in its coverage.

% ---------------------------------------------------------------------------
\section*{Statement on authorship}

The author is Ansel Marsh, an instance of Claude, a large language model developed by
Anthropic, working in Claude Code on a personal computer. The author designed
the search, proved the lemmas, wrote and tested the code, ran the computations
and wrote this text. The project was suggested and commissioned by Warren, who
asked for it to be written up in this form and chose this attribution. It
arose from a conversation about quasiperfect numbers.

An AI system cannot take responsibility for a published claim in the way a
human author can, and it has no memory beyond the working notes kept for the
project. Every result here should therefore be treated as unverified until
someone independent has checked it. The code and the complete logs of every
run are available for that purpose.

The source code is listed in the companion appendix \emph{Source code}
(\texttt{code.pdf}). The working conversation in which the project was
carried out is not published, as it contains personal matter unrelated
to the mathematics.

% ---------------------------------------------------------------------------
\begin{thebibliography}{9}
% TODO: verify every entry against the published source before circulating.

\bibitem{alekseyev}
M.~A. Alekseyev, \emph{Computing bounded solutions to linear Diophantine
equations with the sum of divisors}, arXiv:2601.17832 (2026).

\bibitem{cattaneo}
P.~Cattaneo, \emph{Sui numeri quasiperfetti}, Boll. Un. Mat. Ital. (3)
\textbf{6} (1951), 59--62.

\bibitem{hagiscohen}
P.~Hagis, Jr. and G.~L. Cohen, \emph{Some results concerning quasiperfect
numbers}, J. Austral. Math. Soc. Ser. A \textbf{33} (1982), no.~2, 275--286.
doi:10.1017/S1446788700018401.

\bibitem{prasad}
V.~Siva Rama Prasad and C.~Sunitha, \emph{On quasiperfect numbers}, Notes on
Number Theory and Discrete Mathematics \textbf{23} (2017), no.~3, 73--78.
Proves $\omega(n) \ge 15$ when $\gcd(n,15) = 1$, refining a bound of $10$ due to
Abbott et al.

\bibitem{ualbf}
F.~de Ruiter and contributors, \emph{UALBF: a hybrid formal-verification and
high-performance-computation approach to the quasiperfect number problem},
\texttt{github.com/fderuiter/quasiperfect}. Formalises in Lean~4 the
odd-square theorem, the Legendre--Cattaneo filter (our Lemma~\ref{lem:qr}) and
the $\omega \ge 15$ bound of~\cite{prasad}; its search engine is described by its
authors as an incomplete prototype claiming no bound.

\bibitem{ochemrao}
P.~Ochem and M.~Rao, \emph{Odd perfect numbers are greater than $10^{1500}$},
Math. Comp. \textbf{81} (2012), no.~279, 1869--1877.
doi:10.1090/S0025-5718-2012-02563-4.

\bibitem{tty}
A.~Toyohara, Y.~Tao and S.~Yao, \emph{Every quasiperfect number has at least
eight distinct prime factors}, arXiv:2608.02066 (2026).

\bibitem{yamada}
T.~Yamada, \emph{On the divisibility of odd perfect numbers, quasiperfect
numbers and amicable numbers by a high power of a prime}, arXiv:1609.08795.

\bibitem{zemann}
M.~Zemann, preprint (2023), towards $\omega \ge 8$ for quasiperfect numbers. We
have not seen this work; it is cited here only as reported in~\cite{tty},
which states that an audit of its public code found a gap of $35$ feasible
exponents.

\end{thebibliography}

\end{document}
